\documentclass[11pt]{article}
\usepackage[a4paper,margin=28mm]{geometry}
\usepackage[T1]{fontenc}
\usepackage{lmodern,amsmath,amssymb,amsthm,mathtools,microtype,mathrsfs}
\usepackage[hidelinks]{hyperref}
\usepackage{enumitem}
\setlist{nosep}
\providecommand{\tightlist}{\setlength{\itemsep}{0pt}\setlength{\parskip}{0pt}}
\newtheorem{theorem}{Theorem}[section]
\newtheorem{lemma}[theorem]{Lemma}
\newtheorem{proposition}[theorem]{Proposition}
\newtheorem{corollary}[theorem]{Corollary}
\theoremstyle{definition}
\newtheorem{definition}[theorem]{Definition}
\theoremstyle{remark}
\newtheorem{remark}[theorem]{Remark}
\newcommand{\Z}{\mathbb Z}
\newcommand{\Q}{\mathbb Q}
\newcommand{\R}{\mathbb R}
\newcommand{\Sig}{\Sigma}
\newcommand{\CR}{\mathcal C_{\R}}
\newcommand{\ER}{E_{\R}}
\newcommand{\hbin}{h_{\mathrm{bin}}}
\newcommand{\Stab}{\operatorname{Stab}}
\newcommand{\F}{\mathbb F}
\newcommand{\dimH}{\dim_{\mathrm H}}
\newcommand{\dimP}{\dim_{\mathrm P}}
\newcommand{\dimB}{\overline{\dim}_{\mathrm B}}
\newcommand{\supp}{\operatorname{supp}}
\newcommand{\val}{v_3}
\newcommand{\hb}{h_{\mathrm{bin}}}
\newcommand{\Ram}{\mathcal R}
\title{Zero dimension for the real exceptional set\protect\\
and vanishing $3$-adic nesting constants}
\author{Timofei Iuzhakov}
\date{7 November 2026\\Version 3.0.0}
\allowdisplaybreaks
\emergencystretch=2em
\begin{document}
\maketitle
\begin{abstract}
Let $\Sigma\subset\Z_3$ be the ternary Cantor set with digits $0,1$.
We retain the zero Hausdorff and packing dimension theorems for the
$3$-adic and full real exceptional sets of powers of two, and the
vanishing $3$-adic finite-visit nesting constants proved in versions
1 and 2. All their proofs, including the finite causal-container
appendix, are included.

Version 3 adds two independent results. For Haar-almost every
$b\in\Z_3^\times$, simultaneously for every nonzero $\lambda\in\Z_3$,
at most two distinct nonnegative iterates $b^n\lambda$ belong to
$\Sigma$. We prove an exact finite classification of linear ternary
equations with maximal digital branching and use complete collections
of nonlinear descendants to obtain this metric conclusion.
Independently, if $M_H(X)$ counts positive integers $M\le X$ for which
$\Sigma\cap M^{-1}\Sigma$ has positive Hausdorff dimension, then
$M_H(X)\asymp X$. Linearly many such multipliers have no positive
finite Cantor integer witness. The proof connects terminal carry cycles
to rational points of the real middle-thirds set, proves summability
of their reduced denominators prime to $3$, and applies a divisibility
sieve. This determines the order of growth in Lagarias's second
multiplier-counting problem. Neither new result settles $1\notin E_2$,
the exact asymptotic $M_H(X)\sim X/2$, or the sharp order of the
finite-quotient count.
\end{abstract}

\paragraph{Version history and stable citation.}
This is version 3.0.0 of the same manuscript, with the same author,
title and citation key \texttt{iuzhakov2026lagarias}.
Version 1.0.0 proved zero Hausdorff and packing dimension for the
$3$-adic infinite-visit set $E_2$. Version 2.0.0 added uniform
finite-visit compact covers, vanishing nesting constants and the
full real exceptional-set theorem. Their complete proof sections
are retained below, followed by the new Sections~\ref{metric:section}
and~\ref{sec:positive-density}. Earlier versioned releases remain
unchanged; a citation should identify the version and theorem used.

\paragraph{What is new in version 3.}
Theorem~\ref{metric:main} proves $E_b^{(3)}=\{0\}$ for Haar-almost
every unit base, with one exceptional null set for all nonzero
starting points and all triples of visit times.
Theorem~\ref{metric:classification} classifies every maximal-branching
finite digital plane, including all allowed targets.
Theorems~\ref{thm:pd-main} and~\ref{thm:pd-denominators} prove
$M_H(X)\asymp X$, a linear-size difference
$\mathcal M_H\setminus\mathcal M_C$, and a summable tail of reduced
Cantor denominators. The two new proof branches are independent of
each other and of the earlier dimension and container theorems.
Their connections with earlier work and the exact definitions of the
counting functions are stated in their respective sections.

\paragraph{Which questions are answered.}
The earlier dimension assertions concern Lagarias's Conjectures A and B
in the arXiv version of \cite{LagariasArxiv}, numbered 1.4 and 1.7 in
the journal version \cite{Lagarias}. The new linear-order theorem answers
the order-of-growth part of Problem 2 (journal Problem 1.11). It does
not determine a leading density constant. The varying-base theorem is
not identified with a numbered conjecture in Lagarias's original paper.
The separate quotient manuscript \cite{Quotients} retains its own
version and citation; the new counting theorem is first included here.
No priority over the entire external literature is inferred from
internal review or from publication of this version.

\paragraph{The prescribed point remains unresolved.}
For $b=2$, the point $\lambda=1$ already has the visits $1,4,256$
at times $0,2,8$. Thus $b=2$ is exceptional for the new three-visit
theorem. We do not prove finiteness of its visits, and a zero-dimensional
fixed-base exceptional set can still contain that point.

\section{The $3$-adic finite-visit problem}

We use the metric $|x-y|_3$ on $\Z_3$, with $\val(3)=1$ and
$\val(0)=+\infty$. Put
\[
 \Sigma=\left\{\sum_{j\ge0}x_j3^j:x_j\in\{0,1\}\right\},
 \qquad
 F_k=\{x\in\Z_3:\#\{n\ge0:2^nx\in\Sigma\}\ge k\}.
\]
The visits in this definition occur at distinct nonnegative integer
times. Set
\[
 \Gamma_{\mathrm H}=\lim_{k\to\infty}\dimH F_k,
 \qquad
 \Gamma_{\mathrm P}=\lim_{k\to\infty}\dimP F_k.
\]
These limits exist because the sets decrease. The Hausdorff limit is
the dyadic nesting constant considered in the study of the exceptional
set by Lagarias and by Abram, Bolshakov and Lagarias
\cite{Lagarias,ABL}. We retain the arbitrary-visit convention throughout:
in particular $F_1=\bigcup_{n\ge0}2^{-n}\Sigma$, rather than literally
$\Sigma$.

\begin{theorem}\label{thm:main}
For every $\delta>0$ there is an integer $J\ge1$ and a countable family
of compact subsets $K_m\subset\Z_3$ such that
\[
 F_J\subseteq\bigcup_{m\ge1}K_m,
 \qquad \dimB K_m\le\delta\quad(m\ge1).
\]
Consequently
\[
 \boxed{\Gamma_{\mathrm H}=\Gamma_{\mathrm P}=0.}
\]
\end{theorem}

The distinction from an infinite-visit argument is the order of
quantifiers. The cardinality $J$ will be fixed before any visiting times
are given. The exponent spans of the extracted templates need not be
bounded uniformly. They are fixed separately on each member of the
countable cover before its digit length tends to infinity.

The analytic input is the finite causal-container theorem in Appendix A
of the frozen release v1.0.0 of \cite{Core}. The real theorem is stated
in Theorem~\ref{real:thm:main}; its proof begins in
Section~\ref{sec:real:statement} and uses a separate arithmetic extraction. We state exactly the portion
used in Section~\ref{sec:entropy} and reproduce its complete source
unchanged in Appendix~\ref{app:container}. The new arithmetic extraction is proved in
full below. In the $3$-adic finite-visit argument, neither the extraction
from an infinite set in \cite{Core} nor a uniform bound on the spans
of finite templates is used.

\section{Uniform finite extraction: statement and incidence bound}

For a finite set of nonnegative integers $P$ write
\[
 \Delta(P)=\{v-u:u,v\in P,\ u<v\}.
\]
These are sets, so equal distances within one $P$ are counted only once.

\begin{theorem}[Uniform finite templates]\label{thm:templates}
Fix integers $r\ge3$, $K\ge1$, and $q\ge1$. There is a finite integer
$J=J(r,K,q)$ with the following property. Every set of $J$ distinct
powers $2^n$, $n\ge0$, contains $K$ distinct members $\gamma_i$ and
admits rational units
\[
 \alpha,A_1,\ldots,A_{r-2},u_1,\ldots,u_K
 \in\Q\cap\Z_3^\times
\]
such that
\begin{equation}\label{eq:template}
 c_i:=\gamma_i/\alpha
 =1+\sum_{t=1}^{r-2}A_t3^{e_{ti}}+u_i3^{E_i}
 \qquad(1\le i\le K).
\end{equation}
The positions
\[
 P_i=\{0,e_{1i},\ldots,e_{r-2,i},E_i\}
\]
are in increasing order, every successive gap is at least $q$, and
$\Delta(P_i)\cap\Delta(P_j)=\varnothing$ for $i\ne j$.
The bound $J$ is independent of the original powers and of all exponent
spans. The coefficient $\alpha$ is a virtual unit anchor: it need not
be one of the original powers.
\end{theorem}

We prove this theorem by induction on the number of common monomial
coordinates. At a stage with $d\ge1$ coordinates the data are
\begin{align}
 &\gamma_i=2^{n_i}\quad\text{distinct},\qquad
 0=e_{0i}<e_{1i}<\cdots<e_{d-1,i},\label{eq:stage}\\
 &\gamma_i=\sum_{t=0}^{d-1}a_t3^{e_{ti}}+\varepsilon_i,
 \quad a_t\in\Q\cap\Z_3^\times,
 \quad \val(\varepsilon_i)\ge e_{d-1,i}+q.\label{eq:error}
\end{align}
Successive gaps in \eqref{eq:stage} are at least $q$, and every column
$e_{ti}-e_{si}$, $s<t$, is injective as $i$ varies. For $d=1$ the last
condition is empty.

We use the uniform theorem of Evertse, Schlickewei and Schmidt
\cite[Theorem 1.1]{ESS}. In characteristic zero, an equation
$a_1z_1+\cdots+a_hz_h=1$, with nonzero coefficients and variables in a
subgroup of $(L^\times)^h$ of rank $\rho$, has at most
\[
 \exp\bigl((6h)^{3h}(\rho+1)\bigr)
\]
nondegenerate solutions. Here nondegenerate means that no nonempty
subsum on the left vanishes. The theorem applies over an algebraic
closure of $\Q_3$ and hence to the solutions considered here.

\begin{lemma}[Uniform incidence bound]\label{lem:incidence}
Suppose the $\gamma_i$ are distinct powers of two and every difference
column $e_{ti}-e_{si}$ is injective. For any fixed coefficients
$c_0,\ldots,c_{d-1}\in\Q_3$, the equation
\begin{equation}\label{eq:span}
 \gamma_i=\sum_{t=0}^{d-1}c_t3^{e_{ti}}
\end{equation}
holds for at most
\[
 B_d=\binom d2+\left\lceil
 \exp\bigl((6d)^{3d}(d+2)\bigr)\right\rceil
\]
indices $i$.
\end{lemma}
\begin{proof}
Delete zero coefficients. If no coefficient remains, there is no
solution. For each pair of remaining terms, equality of valuations fixes
one value of $e_{ti}-e_{si}$ and occurs on at most one row. Excluding
at most $\binom d2$ rows, all terms on the right of \eqref{eq:span}
have different valuations; no nonempty subsum of them can vanish.

Choose one remaining index $t_0$ and divide by $c_{t_0}3^{e_{t_0i}}$:
\begin{equation}\label{eq:unit}
 c_{t_0}^{-1}\gamma_i3^{-e_{t_0i}}
 -\sum_{t\ne t_0}\frac{c_t}{c_{t_0}}3^{e_{ti}-e_{t_0i}}=1.
\end{equation}
A vanishing subsum not containing the $\gamma_i$ term is impossible by
the distinct valuations. A vanishing subsum containing that term would,
on using \eqref{eq:span}, force a nonempty complementary subsum of the
original right-hand side to vanish; the complement includes the
$t_0$ term. Thus \eqref{eq:unit} is nondegenerate.

Its variable tuple belongs to the image of
$\langle2\rangle\times\langle3\rangle^d$, a group of rank at most
$d+1$. Two rows with the same tuple have
$2^{n_i-n_j}=3^{e_{t_0i}-e_{t_0j}}$, hence $n_i=n_j$ by multiplicative
independence. Thus rows give distinct solutions. The uniform
unit-equation bound gives the stated number. If only one coefficient
remains, multiplicative independence gives at most one row directly.
\end{proof}

\section{The extension lemma}

\begin{lemma}[One additional coordinate]\label{lem:extension}
For every $d\ge1$ and $L\ge1$ there is a finite $N_d(L)$ such that,
from $N_d(L)$ rows satisfying \eqref{eq:stage}--\eqref{eq:error} and the
injectivity conditions above, one can retain $L$ rows and find rational
units $b_0,\ldots,b_{d-1}$ with
\begin{equation}\label{eq:newresidual}
 \gamma_i=\sum_{t=0}^{d-1}b_t3^{e_{ti}}+u_i3^{E_i},
 \quad u_i\in\Q\cap\Z_3^\times,
 \quad E_i\ge e_{d-1,i}+q,
\end{equation}
such that every column $E_i-e_{ti}$, $0\le t<d$, is injective.
The previous exponent columns do not change. The threshold $N_d(L)$
can be chosen independently of $q$, the coefficients and the exponents.
\end{lemma}

\subsection{Growth order and valuation-minimal minors}
Give the rows any linear order. Color a pair $i<j$ by the signs of all
differences
\[
 (e_{tj}-e_{sj})-(e_{ti}-e_{si})\qquad(s<t).
\]
None is zero. Finite Ramsey theory gives an arbitrarily large
homogeneous subset, provided the original set is sufficiently large.
The numbers
$\Delta_t(i,j)=e_{tj}-e_{ti}$ have the same strict order for every
pair $i<j$: their pairwise comparisons are fixed, and they are
transitive since they compare actual numbers. Temporarily relabel the
coordinates $1,\ldots,d$ in this growth order, so that
\begin{equation}\label{eq:growth}
 \Delta_1(i,j)<\cdots<\Delta_d(i,j)\qquad(i<j).
\end{equation}
This temporary order need not agree with the increasing order of
exponents within a row.

For $s<t$ and $i<j$, \eqref{eq:growth} gives
\begin{equation}\label{eq:monge}
 e_{si}+e_{tj}>e_{sj}+e_{ti}.
\end{equation}
It follows that every square submatrix of $(3^{e_{ti}})$ has exactly
one determinant term of least valuation. Indeed, exchanging an
assignment which is not order-reversing strictly lowers its valuation
by \eqref{eq:monge}. The unique minimum pairs increasing row indices
with decreasing growth-coordinate indices. In particular every
$d$-by-$d$ coordinate minor is nonzero, and its valuation equals that
minimum permutation cost.

Adjoin $\gamma_i$ to the coordinate column and write
\[
 v_i=(3^{e_{1i}},\ldots,3^{e_{di}},\gamma_i)^{\mathsf T}
 \in\Q_3^{d+1}.
\]
Color each $(d+1)$-subset according to whether its augmented determinant
is zero. An all-zero subset of size greater than $B_d$ is impossible:
its augmented matrix would have rank at most $d$, while the coordinate
matrix has rank $d$, forcing the $\gamma_i$ row to be one fixed linear
combination of the coordinate rows, contrary to
Lemma~\ref{lem:incidence}. A second Ramsey refinement therefore gives
as many rows as needed with every augmented $(d+1)$-minor nonzero.

\subsection{Canonical determinant valuations}
For an increasing $(d+1)$-tuple $I=(i_1,\ldots,i_{d+1})$ put
\[
 D_I=\det(v_{i_1},\ldots,v_{i_{d+1}}),\qquad w_I=\val(D_I).
\]
Consider an increasing $(d+2)$-tuple. For each slot
$h\in\{1,\ldots,d+1\}$, compare its two $(d+1)$-subtuples obtained
by keeping either of the adjacent indices in positions $h,h+1$ and
keeping every other index. Color this comparison $t\in\{1,\ldots,d\}$
if the change in $w$ equals the change in $e_t$ at the variable slot;
otherwise color it $G$. At most one coordinate color is possible by
\eqref{eq:growth}. There are at most $(d+1)^{d+1}$ combined colors.
Another finite Ramsey refinement makes the comparison color at each
slot independent of the chosen indices.

If one slot has color $G$, fix the $d$ other indices as pivots, with a
large consecutive group of variable indices in that slot. For these
variables, each column
\begin{equation}\label{eq:goodslot}
 w_{\text{pivots},i}-e_{ti}\qquad(1\le t\le d)
\end{equation}
is injective. This is the situation needed for interpolation below.
We now prove that a $G$ slot must exist on a sufficiently large
homogeneous subset.

Suppose every slot has a coordinate color $t_h$. The graph of increasing
$(d+1)$-tuples, with an edge for a change in only one coordinate, is
connected. Along such an edge the value
$w_I-\sum_h e_{t_h,i_h}$ is unchanged, because the union of the two
tuples is an increasing $(d+2)$-tuple of the type just colored.
Consequently there is one constant $C$ with
\begin{equation}\label{eq:word}
 w_{i_1,\ldots,i_{d+1}}=C+\sum_{h=1}^{d+1}e_{t_h,i_h}
\end{equation}
for every increasing tuple in the homogeneous subset.

For a coordinate $t$ and an increasing $(d+2)$-tuple $I$, the cofactor
identity is
\begin{equation}\label{eq:cofactor}
 \sum_{l=1}^{d+2}(-1)^l3^{e_{t,i_l}}D_{I\setminus i_l}=0.
\end{equation}
Its nonzero summands cannot have a unique least valuation. Color $I$
by a selected pair $(l_t,m_t)$, $l_t<m_t$, of minimum-valuation
positions for each $t$. A further Ramsey refinement fixes these pairs
for all tuples. This uses at most $\binom{d+2}{2}^{d}$ colors.

Write $A_l$ for the valuation of the $l$th summand of
\eqref{eq:cofactor}. For a fixed $t$ and its fixed pair $l<m$, the
identity $A_m-A_l=0$, after inserting \eqref{eq:word}, has the
following slot contributions:
\begin{equation}\label{eq:slotfunctions}
 \begin{cases}
 e_{t_l,i_l}-e_{t,i_l},&\text{at slot }l,\\
 e_{t_j,i_j}-e_{t_{j-1},i_j},&l<j<m,\\
 e_{t,i_m}-e_{t_{m-1},i_m},&\text{at slot }m.
 \end{cases}
\end{equation}
All other contributions cancel. Hold the other slots fixed and give
one slot two possible indices. This is possible for every slot in an
ordered set with at least $d+3$ elements. A difference of two distinct
coordinate functions is injective by hypothesis, so it cannot have
the same value for these two choices. Hence all the labels in each
slot function of \eqref{eq:slotfunctions} coincide, and
\begin{equation}\label{eq:intervalword}
 t_l=t_{l+1}=\cdots=t_{m-1}=t.
\end{equation}
In particular each of the $d$ coordinate labels appears in the word
$t_1,\ldots,t_{d+1}$. Exactly one label appears twice and all the
others appear once.

This word is nonincreasing in the growth order. Indeed, adjacent
summand valuations for coordinate $t$ satisfy
\begin{equation}\label{eq:adjacent}
 A_{j+1}-A_j=
 \Delta_t(i_j,i_{j+1})-\Delta_{t_j}(i_j,i_{j+1}).
\end{equation}
Suppose $t_h=a<b=t_{h+1}$. If $a$ appeared only once, its fixed minimum
pair, by \eqref{eq:intervalword}, would be $h,h+1$. But the next
difference in \eqref{eq:adjacent} is negative, making $A_{h+2}$ smaller
than that minimum. Thus $a$ must be the repeated label. If $b$ appeared
only once, its minimum pair would be $h+1,h+2$, while the preceding
difference in \eqref{eq:adjacent} is positive, making $A_h$ smaller
than that minimum. Thus $b$ must also repeat, a contradiction.
Therefore the word has the form
\begin{equation}\label{eq:descending}
 d,d-1,\ldots,t_*,t_*,\ldots,1.
\end{equation}

\subsection{A two-dimensional residue-field obstruction}
Let the repeated label $t_*$ occupy slots $h,h+1$. Fix $h-1$ pivot
indices before a group of five variable indices, and $d-h$ pivots
after it. There are $d-1$ pivots in total, so $d+4$ available indices
suffice. Their augmented columns are independent, as follows by
extending them to a nonzero $d$-by-$d$ coordinate minor.
Quotient $\Q_3^{d+1}$ by their span. This is a two-dimensional vector
space; let $u_i$ denote the images of the five variable columns.

In any fixed coordinates on the quotient, its pair determinants are
the corresponding augmented determinants times one fixed nonzero
factor (up to signs). By \eqref{eq:word} and \eqref{eq:descending},
\[
 \val\det(u_i,u_j)=C'+e_{t_*,i}+e_{t_*,j}.
\]
The vectors $z_i=3^{-e_{t_*,i}}u_i$ therefore have all pair
determinants nonzero and of the same valuation. There can be at most
four such vectors. To see this without any normalization assumption on
their norms, use $z_1,z_2$ as a basis and write
$z_i=a_i z_1+b_i z_2$ for $i\ge3$. Comparison with the first two
vectors gives $a_i,b_i\in\Z_3^\times$. Comparison between $i,j\ge3$
gives $a_i b_j-a_j b_i\in\Z_3^\times$. Thus the residues of $a_i/b_i$
are distinct nonzero elements of $\F_3$, of which there are only two.
There are at most two vectors beyond the basis, contradicting the
five variables. The assumption that all slots have coordinate colors
is impossible. A $G$ slot exists.

\subsection{Interpolation and preservation of the gap}
Restore the original within-row exponent order
$0=e_{0i}<\cdots<e_{d-1,i}$. Fix the $d$ pivots of a $G$ slot.
Their coordinate matrix is invertible. Let $b_t$ interpolate the
$\gamma_i$ on the pivots, and put $\delta_t=b_t-a_t$. The data for
$\delta_t$ are the errors $\varepsilon_i$.

By Cramer's rule, $\val(\delta_t)\ge q$ for every $t$. In fact, the
denominator determinant has valuation exactly its minimum permutation
cost, by \eqref{eq:monge}. A term of a numerator determinant replaces
one entry $3^{e_{ti}}$ in the corresponding denominator term by
$\varepsilon_i$, whose valuation is at least
$\max_s e_{si}+q\ge e_{ti}+q$. Every numerator term has valuation at
least the denominator's minimum cost plus $q$. Cancellation can only
increase this valuation. The $b_t$ remain units; they are rational,
since both the coordinate matrix and its interpolation data are
rational.

For a variable index put
\[
 R_i=\gamma_i-\sum_{t=0}^{d-1}b_t3^{e_{ti}},\qquad E_i=\val(R_i).
\]
The augmented determinant formula for an interpolation residual gives
$E_i=w_{\text{pivots},i}-C_0$ for one fixed constant $C_0$.
All residuals are nonzero, and \eqref{eq:goodslot} shows that every
$E_i-e_{ti}$ is injective.

It remains to discard rows on which $E_i<e_{d-1,i}+q$. On such a row
the old $\varepsilon_i$ cannot attain the least valuation in
\[
 R_i=\varepsilon_i-\sum_t\delta_t3^{e_{ti}}.
\]
If two \emph{nonzero} error terms tie for the finite least valuation,
that equality fixes one value of the corresponding difference column;
there is at most one such row per unordered pair of coordinates.
These rows number at most $\binom d2$. On every other bad row exactly
one error term, say $t$, is least, giving
\[
 E_i-e_{ti}=\val(\delta_t).
\]
Injectivity of this residual column permits at most one such row for
each $t$, hence at most $d$ more. Deleting at most
$d+\binom d2$ rows gives \eqref{eq:newresidual} and proves
Lemma~\ref{lem:extension}, subject only to choosing the finite
cardinality thresholds, made explicit next.

\subsection{Uniform cardinality bookkeeping}\label{subsec:ramsey}
Let $\Ram_k(m;c)$ be any finite Ramsey number for $c$ colors on
$k$-subsets and monochromatic size $m$; use the usual one-color
convention. For desired output $L$, put
\begin{align*}
 b_d&=d+\binom d2,\qquad M=L+b_d,\\
 H&=\max\left\{d+M,\ \Ram_{d+2}
       \left(d+4;\binom{d+2}{2}^{d}\right)\right\},\\
 P&=\Ram_{d+2}\left(H;(d+1)^{d+1}\right),\\
 Q_0&=\Ram_{d+1}\left(\max\{P,B_d+1\};2\right),\\
 N_d(L)&=\Ram_2\left(Q_0;2^{\binom d2}\right).
\end{align*}
The refinements are respectively the growth comparisons, the nonzero
augmented minors, and the fixed slot colors. A $G$ slot leaves at
least $H-d\ge L+b_d$ variable indices: take the required pivots at
the two ends and all intermediate indices as variables. Discarding
$b_d$ bad rows leaves $L$. If there is no $G$ slot, the minimum-pair
refinement leaves $d+4$ indices, sufficient for every local slot
variation and the five-variable quotient contradiction. This supplies
a finite threshold depending only on $d,L$. In particular, no valuation
depth or exponent span occurs in it.

\section{Iteration and disjoint distance sets}

\begin{proof}[Proof of Theorem~\ref{thm:templates}]
There are $U_q=2\cdot3^{q-1}$ unit residue classes modulo $3^q$.
Pigeonhole sufficiently many original powers into one such class and
choose an ordinary unit representative $a_0$. Then
$\gamma_i=a_0+\varepsilon_i$, with $\val(\varepsilon_i)\ge q$,
provides the $d=1$ starting stage.

Apply Lemma~\ref{lem:extension}. If another common coordinate is
needed, write its residual $R_i=u_i3^{E_i}$ and retain one residue
class of the units $u_i$ modulo $3^q$, represented by a unit $a_d$.
Adjoin the common monomial $a_d3^{e_{di}}$ with $e_{di}=E_i$.
The remaining error has valuation at least $e_{di}+q$.
All former difference columns are unchanged, and all new ones are
injective by the extension lemma. Thus the induction hypotheses hold
with $d+1$ common coordinates.

At $d=r-1$ perform one final extension, leaving its terminal unit
coefficient free. There are then $r$ positions, every successive gap
is at least $q$, and each of the $c=\binom r2$ positive difference
columns is injective. Form the graph on rows in which two rows are
adjacent if their distance sets intersect. A fixed row has at most
$c$ distances. For each of these distances, each of the $c$ difference
columns attains it in at most one row. Thus the graph has maximum
degree at most $c^2$. From $K(c^2+1)$ rows a greedy independent-set
selection retains $K$ pairwise disjoint distance sets. This argument
allows repeated distances within an individual template.

For clarity, all input sizes may be fixed backwards. Start with
$L_{r-1}=K(c^2+1)$ and require $N_{r-1}(L_{r-1})$ rows at that stage.
For each preceding extension choose an output size $U_q$ times the
next stage's required input size, and apply the corresponding $N_d$.
Finally multiply the $d=1$ required input size by $U_q$ for the initial
residue pigeonhole. This is a finite integer $J(r,K,q)$ depending only
on the displayed parameters.

The final common coefficient at exponent zero is a rational unit
$\alpha$. Divide the identity by $\alpha$ to obtain
\eqref{eq:template}. All coefficients and terminal residuals are
rational, because every interpolation used rational data. No bound on
their heights or on the exponent spans has been asserted.
\end{proof}

\section{A finite entropy estimate for weakly separated templates}
\label{sec:entropy}

For $z\in\Z_3$ let $[z]_j\in\{0,\ldots,3^j-1\}$ denote its
residue modulo $3^j$, and define the length-$q$ window
\[
 B_j^{(q)}(z)=\frac{[z]_{j+q}-[z]_j}{3^j}.
\]
Let $\Sigma_L=\{\sum_{j=0}^{L-1}x_j3^j:x_j\in\{0,1\}\}$.
All entropies and logarithms below use the natural logarithm.

We first specify the external finite result. This is Appendix A,
``Exact theorem'', of \cite{Core}.

\begin{theorem}[Finite causal-container input]\label{thm:container}
Let $r\ge3$, $V_0=\{0,\ldots,N-1\}$, and $\Omega$ be an alphabet
of size $Q$. A nonempty finite family $\mathcal W\subseteq\mathcal X^L$
determines letters $Z_j(x)\in\Omega$; these maps need not be local.
Let a fixed family of ordered positional edges
$e=(j_0<\cdots<j_{r-1})$ have relations
$\mathcal R_e(x_{<j_0})\subseteq\Omega^r$ depending only on the
indicated input prefix, and suppose every word satisfies every relation.
Assume that distinct edges with the same source have disjoint future
position sets and that each unordered pair of positions belongs to
at most $C$ edges.

For positive integer thresholds $T_{\mathrm s},T_2,\ldots,T_{r-1}$
put $D=CQ^{r-3}$ and
\begin{equation}\label{eq:theta}
 \theta=Q\left[
 \frac{D\sum_{j=2}^{r-1}T_j+1}{T_{\mathrm s}}
 +\frac1{T_2}
 +D\sum_{j=3}^{r-1}\frac{1+\sum_{l=2}^{j-1}T_l}{T_j}
 \right].
\end{equation}
If $\theta\le1/2$, each word has a fingerprint $F(x)$ and nonempty
assigned alphabets $A_j(x)\subseteq\Omega$ such that $Z_j(x)\in A_j(x)$,
$A_j(x)$ is determined by $(x_{<j},F(x))$, and, outside at most
\begin{equation}\label{eq:badbound}
 Q\left(T_{\mathrm s}+D\sum_{j=2}^{r-1}T_j\right)N
\end{equation}
edges, the \emph{full Cartesian product} of their assigned alphabets
is contained in $\mathcal R_e(x_{<j_0})$. For every distribution on
$\mathcal W$,
\[
 H(F)\le N\bigl[\hb(\theta)+\theta\log Q\bigr],
 \quad \hb(t)=-t\log t-(1-t)\log(1-t).
\]
\end{theorem}

The proof of Theorem~\ref{thm:container} is reproduced in
Appendix~\ref{app:container} from the cited frozen source. The remaining argument, including its application to
unrestricted ternary inputs and rational units, is given here.

\begin{proposition}[Finite-word bound]\label{prop:words}
Choose $b\ge1$ so that $3^b>2(r-1)$ and put $q=b+1$.
Suppose \eqref{eq:template} has successive position gaps at least
this $q$ and pairwise disjoint positive distance sets, as in
Theorem~\ref{thm:templates}. Set
\[
 q=b+1,\quad Q=3^{(r-1)q},\quad
 C=\binom r2,\quad D=CQ^{r-3},\quad R=\max_i E_i.
\]
For $L\ge q$ let $W_L$ be the set of all ternary input words modulo
$3^L$ for which
$[c_iX]_L\in\Sigma_L$ for every $1\le i\le K$; no infinite
extendability is required. Set $N=L-q+1=L-b$. With thresholds as in
Theorem~\ref{thm:container}, define
\[
 \epsilon_F=\hb(\theta)+\theta\log Q,
 \qquad \eta=\frac{Q(T_{\mathrm s}+D\sum_{j=2}^{r-1}T_j)}{K}.
\]
If $\theta\le1/2$, then
\begin{equation}\label{eq:finitebound}
 \log|W_L|\le
 N\left[\epsilon_F+\frac{\log2}{r}+\eta\log3\right]
 +(R+b)\log3.
\end{equation}
\end{proposition}

\subsection{Ordinary carries and the positional graph}
For a source $j$ and a template $i$ with $v=j+E_i<N$, consider the
$r$ window positions
\[
 j,\quad v-e_{r-2,i},\ldots,v-e_{1i},\quad v.
\]
At these positions use respectively the windows of $X$, the
corresponding $A_tX$, and $X$ again. The joint letter available at
every position is
\begin{equation}\label{eq:joint}
 Z_j(X)=\bigl(B_j^{(q)}(X),
 B_j^{(q)}(A_1X),\ldots,B_j^{(q)}(A_{r-2}X)\bigr).
\end{equation}
It takes values in an alphabet of size $Q$.

At position $v$, the identity for $c_iX$ has $r$ summands:
$X$, the $r-2$ terms $3^{e_{ti}}A_tX$, and $3^{E_i}u_iX$.
Their ordinary nonnegative residues below position $v$ sum to a
carry $h\in\{0,\ldots,r-1\}$ into that position. The terminal summand
has window, modulo $3^q$,
\[
 u_i B_j^{(q)}(X)+D_{ij}(X_{<j}),\qquad
 D_{ij}=\frac{3^{E_i}u_i[X]_j-[3^{E_i}u_i[X]_j]_v}{3^v}.
\]
Here $D_{ij}\in\Z_3$ depends only on the input digits before $j$.
Consequently the corresponding edge relation is
\begin{equation}\label{eq:carry}
 u_iB_j^{(q)}(X)+\sum_{t=1}^{r-2}B_{v-e_{ti}}^{(q)}(A_tX)
 +B_v^{(q)}(X)+D_{ij}+h\in\Sigma_q\pmod{3^q}
\end{equation}
for some $0\le h\le r-1$. This derivation is valid for rational
$3$-adic units, with all residues interpreted in $\Z_3$.
It imposes no restriction $X\in\Sigma$.

Every future displacement from $j$ is a positive distance in $P_i$.
Thus different templates have disjoint future sets at the same source.
If two positions belong to an edge, their distance belongs to exactly
one of the template distance sets, and there are at most
$\binom r2$ choices of roles in that template. The source is then
fixed. The graph therefore satisfies the two hypotheses of
Theorem~\ref{thm:container} with the stated $C$. The relations
\eqref{eq:carry} are causal, as their only extra parameter is $D_{ij}$.

\subsection{The one-digit Cartesian-product bound}
Freeze the low $b=q-1$ digits of each of the $r$ windows. Let each
remaining high digit range in a nonempty subset of $\F_3$. Suppose
the full Cartesian product of these choices satisfies
\eqref{eq:carry}, allowing the target and carry to vary with the choice.
For a target $y\in\Sigma_q$, its low part satisfies
$0\le[y]_b\le R_b=(3^b-1)/2$. Hence
\[
 [y]_b-h\in[-(r-1),R_b].
\]
This interval has length $R_b+r-1<3^b$. For the fixed residue modulo
$3^b$ determined by the frozen inputs and $D_{ij}$, it therefore
contains at most one representative. After subtracting that
representative and dividing by $3^b$, the sum of the high-digit factors,
with the first multiplied by the unit $u_i\bmod3$, lies in a single
translate of $\{0,1\}\subset\F_3$.

No factor can have size three. Two factors of size at least two would
already have sumset $\F_3$. Thus at most one factor has size two and
all the others have size one. The product of the factor cardinalities
is at most two, or equivalently their log-cardinalities sum to at most
$\log2$.

\subsection{Conditional supports, incidence, and entropy}
Apply Theorem~\ref{thm:container} to the uniform distribution on $W_L$,
using \eqref{eq:joint} and \eqref{eq:carry}. (The set is nonempty
because it contains zero.) For each realized word and $0\le j<N$ put
\[
 C_j=\supp(X_{j+b}\mid X_{<j+b},F),\qquad w_j=\log|C_j|\le\log3.
\]
Under this conditioning the low $b$ digits of every window of $X$ or
$A_tX$ starting at $j$ are fixed. Multiplication by a unit maps the
high input digit bijectively onto the high output digit, so the
projected support of the appropriate window has cardinality $|C_j|$.
Moreover it lies in the relevant projection of the assigned alphabet
$A_j$: every word in this finer conditional atom has the same
$(X_{<j},F)$ as the realized word.

On a good edge the full-product conclusion of
Theorem~\ref{thm:container} permits these projected supports to be
combined, although the supports at different positions come from
different prefix atoms. Freezing the low digits to their realized
values and using the preceding Cartesian-product bound gives
\[
 \sum_{j\in e}w_j\le\log2.
\]
A bad edge instead costs at most $r\log3$. Let $E_{\mathrm{bad}}$
be the number of bad edges for the realized word. There are at most
$KN$ edges. For a given template, each role covers all positions
except at most $E_i\le R$ boundary positions. Summing incidences
over all roles and templates gives
\[
 rK\sum_{j=0}^{N-1}w_j-rKR\log3
 \le\sum_e\sum_{j\in e}w_j
 \le KN\log2+rE_{\mathrm{bad}}\log3.
\]
The estimate remains valid when some templates have no edges because
$E_i\ge N$, since the stated boundary loss then dominates the
missing incidences. Dividing by $rK$ and using
$E_{\mathrm{bad}}\le\eta KN$ yields
\[
 \sum_jw_j\le N\frac{\log2}{r}+\eta N\log3+R\log3.
\]
The windows' high input digits are precisely the positions
$b,b+1,\ldots,L-1$. The first $b$ digits cost at most $b\log3$.
The chain rule and the support bound on conditional entropy give
\[
 H(X\mid F)\le b\log3+\mathbb E\sum_jw_j.
\]
Finally $H(X)\le H(F)+H(X\mid F)$, and the fingerprint bound is
$H(F)\le N\epsilon_F$. Since $X$ is uniform on $W_L$, this proves
\eqref{eq:finitebound} and Proposition~\ref{prop:words}.

\section{Compact covers and the limit of finite-visit dimensions}

Fix $r\ge3$. Choose $b$, $q=b+1$, $Q,C,D$ as above. For an integer
$u\ge2$ set
\begin{equation}\label{eq:thresholds}
 T_j=u^{2^{j-1}-1}\ (2\le j\le r-1),\qquad
 T_{\mathrm s}=u^{2^{r-1}-1},\qquad K=u^{2^r}.
\end{equation}
Every ratio occurring in \eqref{eq:theta} tends to zero as
$u\to\infty$. Also
$Q(T_{\mathrm s}+D\sum_jT_j)/K\to0$. Thus
\begin{equation}\label{eq:vanishing}
 \theta\longrightarrow0,
 \qquad \epsilon_F\longrightarrow0,
 \qquad \eta\longrightarrow0
 \quad\text{for fixed }r.
\end{equation}

For any one extracted finite template system $\mathcal D$ define
\[
 K_{\mathcal D}=\{X\in\Z_3:c_iX\in\Sigma\text{ for }1\le i\le K\}.
\]
It is compact. Its number of occupied length-$L$ cylinders is at most
$|W_L|$. Proposition~\ref{prop:words}, applied with the fixed span
$R$ of this system and then $L\to\infty$, gives
\begin{equation}\label{eq:box}
 \dimB K_{\mathcal D}\le
 s_{r,u}:=\frac{\log2}{r\log3}+\frac{\epsilon_F}{\log3}+\eta.
\end{equation}
The fact that $W_L$ contains all locally valid finite words only makes
this upper bound larger; an assumption of infinite extendability is
unnecessary.

\begin{proof}[Proof of Theorem~\ref{thm:main}]
Given $\delta>0$, first choose $r$ so that
$\log2/(r\log3)<\delta/2$. Then choose a finite $u$ sufficiently
large that $\theta\le1/2$ and
$\epsilon_F/\log3+\eta<\delta/2$. Fix the associated $K,q$ and
$J=J(r,K,q)$ from Theorem~\ref{thm:templates}.

For any $x\in F_J$, choose $J$ distinct true visit slopes
$\gamma=2^n$ and apply that finite theorem. The normalized variable
$X=\alpha x$ lies in $K_{\mathcal D}$, since
$c_iX=\gamma_i x\in\Sigma$. The normalized input $X$ need not lie
in $\Sigma$, which is why Proposition~\ref{prop:words} was stated
for arbitrary ternary inputs.

All coefficients constructed in the extraction are rational, and all
exponents and visit times are ordinary integers. There are only
countably many possible finite systems $\mathcal D$, including their
unit anchors $\alpha_{\mathcal D}$. Hence
\[
 F_J\subseteq\bigcup_{\mathcal D}
 \alpha_{\mathcal D}^{-1}K_{\mathcal D}.
\]
These are compact sets, and multiplication by a $3$-adic unit is an
isometry. By \eqref{eq:box}, each has upper box dimension at most
$s_{r,u}<\delta$. This proves the covering statement.

Packing dimension is countably stable and is bounded above by upper
box dimension on bounded sets. Therefore
\[
 \dimH F_J\le\dimP F_J\le\delta.
\]
As $F_{k+1}\subseteq F_k$ and $\delta>0$ was arbitrary, both limiting
dimensions are zero.
\end{proof}

\begin{remark}[Fixed finite-rank groups]
The same proof applies to any fixed subgroup
$G\le\Q^\times\cap\Z_3^\times$ of finite rank $\rho$. Replace
distinct powers of two by distinct members $\gamma_i\in G$ and put
$F_k(G)=\{x\in\Z_3:\#\{\gamma\in G:\gamma x\in\Sigma\}\ge k\}$.
In Lemma~\ref{lem:incidence} the variable group has rank at most
$\rho+d$, so replace the factor $d+2$ in $B_d$ by $\rho+d+1$.
The normalized solution tuples are still distinct: equality of their
first coordinates gives $\gamma_i/\gamma_j=3^{e_{t_0i}-e_{t_0j}}$;
the left side is a unit, so the exponent difference is zero and
$\gamma_i=\gamma_j$. All remaining arguments are unchanged.
The coefficients stay rational, and the compact-cover data are
countable. Thus for each $\delta>0$ some finite
$J=J(\delta,\rho)$ gives the same compact-cover conclusion for
$F_J(G)$, and both finite-visit dimension limits are zero.
The finite-rank hypothesis is essential to this argument; this remark
does not concern arbitrary families of integer multipliers.
\end{remark}

\begin{corollary}
The infinite-visit exceptional set
\[
E_2=\{x\in\Z_3:2^nx\in\Sigma\text{ for infinitely many }n\ge0\}
\]
has Hausdorff and packing dimension zero.
\end{corollary}
\begin{proof}
$E_2\subseteq F_J$ for every $J$, so the result follows from
Theorem~\ref{thm:main}.
\end{proof}

The finite-visit conclusion supplies uniformity in the number of
visits, rather than a statement about a particular initial point.
It does not imply that $1\notin E_2$, that every nonzero orbit has
finitely many visits, or that the exceptional set is countable.
In particular no resolution of the pointwise Erd\H{o}s problem is
claimed. The effective constants supplied by repeated Ramsey numbers
and the unit-equation bound are far too large for the result to serve
as a practical search cutoff.

\section{The real exceptional set and the proof dependencies}\label{sec:real:statement}
Define
\begin{align*}
 \CR&=\left\{\sum_{j=-\infty}^{m}\varepsilon_j3^j:
       m\in\Z,\ \varepsilon_j\in\{0,1\}\right\}\subseteq[0,\infty),\\
 \ER&=\{\lambda>0:2^n\lambda\in\CR\text{ for infinitely many }n\ge0\}.
\end{align*}
The digit condition includes every fractional digit. This is the
untruncated real exceptional set in Lagarias~\cite[\S1.2]{Lagarias},
whose Conjecture~1.4 asks whether $\dimH\ER=0$. The corresponding
statement is called Conjecture~A in the arXiv manuscript~\cite{LagariasArxiv}.
It differs from the truncated set defined using only
$\lfloor2^n\lambda\rfloor$: Lagarias~\cite[Theorem~1.3]{Lagarias}
proved that the latter has Hausdorff dimension $\log_3 2$.

\begin{theorem}\label{real:thm:main}
The untruncated real exceptional set satisfies
\[
                    \boxed{\dimH\ER=\dimP\ER=0.}
\]
More precisely, for every $\delta>0$ there is a countable cover of
$\ER$ by compact subsets of $\R$ of upper box dimension at most $\delta$.
\end{theorem}

The real proof uses the same finite causal-container theorem,
Theorem~\ref{thm:container}, whose complete frozen proof from
version~1.0.0 is reproduced in Appendix~\ref{app:container}.
Its two-translate capacity estimate is Lemma~\ref{vone:lem:capacity}
below; we reproduce its proof from \cite[Lemma~2.1]{Core} as well.
Thus neither finite input is left implicit. We give the full affine
finite-word application, the transfer to real prefixes, and the dimension
argument. We also restate and prove the arithmetic extraction for an
arbitrary infinite set of visit times. Its external input is the
finite-rank unit-equation theorem of Evertse, Schlickewei and
Schmidt~\cite{ESS}. This infinite extraction uses a true visiting anchor;
it is distinct from the uniform finite extraction with a virtual anchor
in Theorem~\ref{thm:templates}.

Throughout, logarithms and Shannon entropies are natural. Write
\[
 \Sig_L=\left\{\sum_{j=0}^{L-1}\varepsilon_j3^j:
              \varepsilon_j\in\{0,1\}\right\}
              \subseteq\Z/3^L\Z.
\]
For $z\in\Z_3$, let $[z]_m$ be its representative modulo $3^m$ in
$[0,3^m)$, and put
\[
 B_{j,q}(z)=\left[(z-[z]_j)/3^j\right]_q.
\]
A rational number with denominator prime to $3$ is interpreted in
$\Z_3$ when residues or valuations are used. Its ordinary real value
is used in floors and real intervals. This distinction is essential
in the transfer below.

\section{Templates from an infinite set of exponents}
The following extraction is the version of
\cite[Lemma~3.1]{Core} that separates the arithmetic from the space in
which the initial value lies.

\begin{lemma}\label{real:lem:extraction}
Let $S\subseteq\Z_{\ge0}$ be infinite, and fix integers
$r\ge3$, $K\ge1$ and $q\ge1$. There exist distinct indices
$n_0,n_1,\ldots,n_K\in S$ and ordinary integer units
$A_1,\ldots,A_{r-2}\in\Z_3^\times$ such that
$c_i=2^{n_i-n_0}$ has an exact expansion
\begin{equation}\label{real:eq:template}
 c_i=1+\sum_{t=1}^{r-2}A_t3^{e_{ti}}+u_i3^{E_i},
 \qquad u_i\in\Z_3^\times,
\end{equation}
where $0<e_{1i}<\cdots<e_{r-2,i}<E_i$, every successive gap,
including the gap from $0$, is at least $q$, and
\begin{equation}\label{real:eq:separation}
 \min\{e_{1,i+1},e_{2,i+1}-e_{1,i+1},\ldots,
                 E_{i+1}-e_{r-2,i+1}\}>E_i\qquad(1\le i<K).
\end{equation}
\end{lemma}
\begin{proof}
The distinct slopes $\alpha=2^n$, $n\in S$, lie in the compact group
$\Z_3^\times$. Choose a subsequence $\alpha_j\to\rho\in\Z_3^\times$.
After discarding a possible equality, write
\[
 \rho^{-1}\alpha_j=1+v_{1j}3^{e_{1j}},
 \qquad v_{1j}\in\Z_3^\times,\quad e_{1j}\to\infty.
\]
Pass to a subsequence on which $v_{1j}\to a_1\in\Z_3^\times$.
Whenever $v_{1j}-a_1\ne0$, normalize this difference by its valuation
and take a further convergent subsequence of its unit part. Repeating
this procedure produces successive exponent gaps tending to infinity.

An infinite collection of exact stops after $s$ terms would give
infinitely many solutions of
\[
 \rho^{-1}\gamma-\sum_{t=1}^{s}a_t3^{e_t}=1,
 \qquad \gamma\in\langle2\rangle,\quad 0<e_1<\cdots<e_s.
\]
The coordinate variables belong to the finite-rank multiplicative group
$\langle2\rangle\times\langle3\rangle^s$.
Every nonempty subsum on the left is nonzero: its terms have distinct
3-adic valuations $0,e_1,\ldots,e_s$, so the term of least valuation
cannot cancel. The characteristic-zero finite-rank theorem of
\cite{ESS}, applied over an algebraic closure of $\Q_3$, excludes
infinitely many such nondegenerate solutions. Exact stops may therefore
be discarded at each of the finitely many stages required here.

We obtain
\[
 \rho^{-1}\alpha_j
 =1+\sum_{t=1}^{r-2}a_t3^{e_{tj}}+u_j^*3^{E_j},
 \qquad a_t,u_j^*\in\Z_3^\times,
\]
with every successive gap tending to infinity. Choose $K$ members,
now indexed by $i$, so that the required lower gap bounds and
\eqref{real:eq:separation} hold. Choose ordinary integer units $A_t$ with
\[
 v_3(A_t-a_t)\ge\max_i(E_i-e_{ti}+q).
\]
Finally choose another member $\alpha_0=2^{n_0}$ of the original
subsequence, outside the selected finite set, with
\[
 v_3(\rho/\alpha_0-1)\ge\max_i(E_i+q).
\]
In replacing $\alpha_i/\rho$ by $\alpha_i/\alpha_0$, and $a_t$ by
$A_t$, the total error in row $i$ is divisible by $3^{E_i+q}$ in
$\Z_3$. The normalized final remainder is consequently still a unit.
This gives \eqref{real:eq:template} with the same exponents and separation.
\end{proof}

\begin{remark}\label{real:rem:rational}
The slopes $c_i$ are positive rational 3-adic units, but need not be
ordinary integers: the selected anchor $n_0$ can exceed some $n_i$.
The proof used only the infinite set $S$, not an initial point in
$\Z_3$. In particular, it applies to true visit times of a real point.
\end{remark}

\section{An affine finite-word estimate}\label{sec:real:finite}
Fix a template system as in Lemma~\ref{real:lem:extraction}. Choose
$k\ge1$ and $b\ge1$ with $3^b>r-1$, and use $q=b+k$ in the
extraction. Set
\[
 R=\max_iE_i,\qquad Q=3^{(r-1)q},\qquad
 C=\binom r2,\qquad D=CQ^{r-3}.
\]
For positive integer thresholds $T_{\rm s},T_2,\ldots,T_{r-1}$ define
\begin{align}
 \theta&=Q\left[\frac{D\sum_{j=2}^{r-1}T_j+1}{T_{\rm s}}
 +\frac1{T_2}+D\sum_{j=3}^{r-1}
          \frac{1+\sum_{l=2}^{j-1}T_l}{T_j}\right],\label{real:eq:theta}\\
 \epsilon_F&=\hbin(\theta)+\theta\log Q,
 &\hbin(t)&=-t\log t-(1-t)\log(1-t),\label{real:eq:epsF}\\
 \eta&=\frac{Q(T_{\rm s}+D\sum_{j=2}^{r-1}T_j)}{K},
 &s&=\frac{\epsilon_F}{\log3}+\frac1k+
       \frac{2^{k+1}\log2}{rk\log3}+\eta.\label{real:eq:s}
\end{align}
We assume $\theta\le1/2$.

The following estimate is reproduced from \cite[Lemma~2.1]{Core}.
It is uniform in both target translates.

\begin{lemma}[Two-translate capacity]\label{vone:lem:capacity}
Suppose $A_1,\ldots,A_r$ are nonempty subsets of $\Z/3^k\Z$ and
\[
 T=A_1+\cdots+A_r\subseteq(z_1+\Sigma_k)\cup(z_2+\Sigma_k).
\]
Then
\begin{equation}\label{vone:eq:capacity}
 \sum_{i=1}^r\log|A_i|\le r\log3+2^{k+1}\log2=:B_{r,k}.
\end{equation}
\end{lemma}
\begin{proof}
Let $H=\Stab(T)$ and $|H|=3^h$. A coset of $H$ fixes the lowest
$k-h$ digits. Its intersection with a translate of $\Sigma_k$ has at
most $2^h$ elements: after translating back, the highest $h$ digits
form a word of $\Sigma_h$. Since the nonempty $H$-periodic set $T$
contains a whole coset, $3^h\le2^{h+1}$. Thus $h\le1$.

Set $n_i=|A_i+H|/|H|$. Kneser's addition theorem \cite{DeVos}, applied
in the quotient by $H$, gives
\[
 \sum_i(n_i-1)\le |T|/|H|-1\le 2^{k+1}/|H|-1.
\]
Indeed the final sum in the quotient has trivial stabilizer; the
stabilizer of any partial sum is contained in that of the final sum,
so the two-summand aperiodic inequality may be applied successively.
Since $n\le2^{n-1}$ for positive integers $n$,
\[
 \sum_i\log|A_i|\le rh\log3+\sum_i(n_i-1)\log2
                  \le r\log3+2^{k+1}\log2.
\]
\end{proof}

\begin{proposition}\label{real:prop:finite}
For arbitrary fixed shifts $\beta_1,\ldots,\beta_K\in\Z_3$ and
$L\ge q$, let
\[
 W_L(\boldsymbol\beta)=
 \{x\in\{0,\ldots,3^L-1\}:
      [x]_L\in\Sig_L,\ [c_ix+\beta_i]_L\in\Sig_L\ (1\le i\le K)\}.
\]
Then, with $N=L-q+1$,
\begin{equation}\label{real:eq:finite}
 \log|W_L(\boldsymbol\beta)|
 \le Ns\log3+(R+b+k-1)\log3
\end{equation}
whenever the set is nonempty. Equivalently, including the empty case,
\begin{equation}\label{real:eq:finite-simple}
 |W_L(\boldsymbol\beta)|\le 3^{R+b+k-1}\,3^{sL}.
\end{equation}
Both bounds are uniform in the shifts. All locally admissible finite
words are counted; an extension to an infinite admissible word is not
required.
\end{proposition}
\begin{proof}
We give the affine application in detail because uniformity in the
shifts is needed for the real transfer. View a finite input word as its
ordinary representative $X\in[0,3^L)$, extended by zero digits if needed,
and set $Y_t=A_tX$ in $\Z_3$.

\emph{Exact causal arithmetic.}
For template $i$ and source position $j$, put $v=j+E_i$. The participating
windows are
\[
 B_0=B_{j,q}(X),\quad
 B_t=B_{v-e_{ti},q}(Y_t)\ (1\le t\le r-2),\quad
 B_{\rm end}=B_{v,q}(X).
\]
Write $P=[X]_j$ and define
\[
 T=\beta_i+3^{E_i}u_iP,\qquad
 t_v=[T]_v,\qquad D_j=(T-t_v)/3^v\in\Z_3.
\]
The center $D_j$ depends only on the fixed data and the input prefix
$X_{<j}$. In the exact identity for $c_iX+\beta_i$, the lower parts are
\[
 [X]_v,\qquad 3^{e_{ti}}[Y_t]_{v-e_{ti}}\ (1\le t\le r-2),\qquad t_v.
\]
These are $r$ nonnegative integers below $3^v$, so their ordinary carry
lies in $\{0,\ldots,r-1\}$. Taking the next $q$ digits gives the relation
\begin{equation}\label{real:eq:causal}
 u_iB_0+\sum_{t=1}^{r-2}B_t+B_{\rm end}+D_j+e\in\Sig_q
       \pmod{3^q}\qquad\text{for some }0\le e\le r-1.
\end{equation}
Only windows contained in the first $L$ output digits will be used.
Thus this relation is valid for every word in $W_L(\boldsymbol\beta)$.

\emph{Low digits and capacity.}
On any Cartesian product of windows satisfying \eqref{real:eq:causal},
fix the lowest $b$ digits of each argument. The low residue of their
sum, including $D_j$, is fixed. If $y\in\Sig_q$ is an allowed output,
then
\[
                 y_{\rm low}-e\in[-(r-1),3^b-1].
\]
Since $3^b>r-1$, a fixed residue modulo $3^b$ has at most two
representatives in this interval. Dividing the corresponding identities
by $3^b$ shows that the sum of the high $k$-digit factors lies in at most
two translates of $\Sig_k$. The first high factor is multiplied by the
unit $u_i$, an affine bijection after its low digits are fixed.
We use Lemma~\ref{vone:lem:capacity}: for nonempty
$H_1,\ldots,H_r\subseteq\Z/3^k\Z$,
\begin{equation}\label{real:eq:capacity}
 H_1+\cdots+H_r\subseteq(z_1+\Sig_k)\cup(z_2+\Sig_k)
 \quad\Longrightarrow\quad
 \sum_{\nu=1}^r\log|H_\nu|\le r\log3+2^{k+1}\log2=:B_{r,k}.
\end{equation}
The proof above derives this estimate from Kneser's theorem;
its constants do not depend on either translate.

\emph{The finite causal-container input.}
At each position $j\in\{0,\ldots,N-1\}$ take the joint letter
\[
 Z_j=(B_{j,q}(X),B_{j,q}(Y_1),\ldots,B_{j,q}(Y_{r-2}))
\]
in an alphabet of size $Q$. The ordered positional edges are
\[
 (j,j+E_i-e_{r-2,i},\ldots,j+E_i-e_{1i},j+E_i)
\]
whenever all positions belong to $\{0,\ldots,N-1\}$.
Their relations select the indicated coordinates in \eqref{real:eq:causal}.
By \eqref{real:eq:separation}, future position sets of distinct edges from
the same source are disjoint. Every pair of positions belongs to at
most $C=\binom r2$ edges: the distance belongs to at most one template,
and that template has at most $\binom r2$ role pairs.

Theorem~\ref{thm:container}, with its complete proof in
Appendix~\ref{app:container}, therefore supplies a deterministic
fingerprint $F$ and alphabets $\mathcal A_j$, each determined by
$(X_{<j},F)$ and containing the actual joint letter. It gives
\begin{equation}\label{real:eq:budgets}
 H(F)\le N\epsilon_F,\qquad
 E_{\rm bad}\le Q\left(T_{\rm s}+D\sum_{j=2}^{r-1}T_j\right)N.
\end{equation}
Except for at most $E_{\rm bad}$ edges, the \emph{full Cartesian product}
of the assigned alphabets satisfies its relation at the actual source
prefix. This full-product conclusion, rather than compatibility along
one input, is the property used next.

\emph{Conditional supports and entropy.}
Take $X$ uniformly distributed on the nonempty set
$W_L(\boldsymbol\beta)$. For each realized word put
\[
 C_j=\supp(X_{[j+b,j+b+k)}\mid X_{<j+b},F),\qquad
 w_j=\log|C_j|\le k\log3.
\]
Fixing $X_{<j+b}$ fixes the low $b$ digits of all joint windows at $j$.
Multiplication by any $A_t$ induces an affine bijection on the next
$k$ digits. Each required high-coordinate projection of the conditional
joint support consequently has exactly $|C_j|$ elements. This joint
support is contained in $\mathcal A_j$, since $\mathcal A_j$ is determined
by the shorter prefix $X_{<j}$ and the same fingerprint.

On a good edge these projections may be selected independently, even
though the prefixes defining them differ at the different positions.
The full-product property applies to every such selection; no common
new input realizing all selections is asserted. The low-digit reduction
and \eqref{real:eq:capacity} imply, pointwise on that edge,
\[
                     \sum_{\nu=0}^{r-1}w_{j_\nu}\le B_{r,k}.
\]
On a bad edge use $rk\log3$. Template $i$ has
$\max(N-E_i,0)$ edges, and each of its $r$ roles misses
$\min(E_i,N)$ positions. Accounting for these boundary incidences gives
\[
 rK\sum_{j=0}^{N-1}w_j
 \le B_{r,k}KN+rk\log3\,E_{\rm bad}+rk\log3\,KR,
\]
and hence
\begin{equation}\label{real:eq:weights}
 \sum_jw_j\le \frac{NB_{r,k}}r+
                       k\log3\left(\frac{E_{\rm bad}}K+R\right).
\end{equation}
For $h_t=H(X_t\mid X_{<t},F)$ the entropy chain rule yields
\[
 \sum_{t=j+b}^{j+b+k-1}h_t\le\mathbb E w_j.
\]
The multiplicities of the $h_t$ in this sum are at most $k$ and total
$kN$. Since $h_t\le\log3$, their deficit from $kL$ is at most
$k(q-1)\log3$ in entropy units. Therefore
\[
 H(X\mid F)\le k^{-1}\sum_j\mathbb E w_j+(q-1)\log3.
\]
Combining this with \eqref{real:eq:budgets}--\eqref{real:eq:weights}, and using
$H(X)=H(F)+H(X\mid F)=\log|W_L(\boldsymbol\beta)|$, proves
\eqref{real:eq:finite}. Dropping its nonpositive contribution
$-(q-1)s\log3$ gives \eqref{real:eq:finite-simple}. All bounds are independent
of the shifts, as required.
\end{proof}

\section{Transfer from real truncations to finite words}
The next lemma is the additional step connecting the two number systems.
It requires a bound for all finite words, not merely for prefixes that
extend to infinite admissible 3-adic inputs.

\begin{lemma}[Finite-word transfer]\label{real:lem:transfer}
Let $c_i=p_i/q_i>0$ be rational numbers in lowest terms, with
$3\nmid p_iq_i$, for $1\le i\le K$. Suppose that, for some constants
$C_0>0$, $s\ge0$ and $L_0$, their affine word sets satisfy
\[
 |W_L(\boldsymbol\beta)|\le C_0\,3^{sL}
 \quad(L\ge L_0)
\]
for all $\boldsymbol\beta\in\Z_3^K$. For each integer $M\ge1$, put
\[
 K_{\R}(\boldsymbol c;M)=
 \{X\in[0,M]:X\in\CR,\ c_iX\in\CR\ (1\le i\le K)\}.
\]
This set is compact and has upper box dimension at most $s$.
\end{lemma}
\begin{proof}
On a bounded interval, there are only finitely many possible integer
digit prefixes in $\CR$. Each has a compact set of fractional tails,
the continuous image of the compact product $\{0,1\}^{\mathbb N}$.
It follows that $\CR$ is closed on bounded intervals, and the stated
finite intersection is compact.

Fix $X\in K_{\R}(\boldsymbol c;M)$ and $L\ge1$. Define
\[
 T_L=\lfloor3^LX\rfloor,\qquad
 T_{i,L}=\lfloor3^Lc_iX\rfloor,\qquad
 U_L=T_L\bmod3^L\in\{0,\ldots,3^L-1\}.
\]
The full real digit conditions imply
$U_L\in\Sig_L$ and $T_{i,L}\bmod3^L\in\Sig_L$.
Indeed, after multiplication by $3^L$, the remaining fractional
$0/1$ tail lies in $[0,1/2]$, so taking the floor gives exactly the
allowed integer prefix. This also handles terminating expansions.
Set
\[
 \beta_{i,L}=T_{i,L}-c_iT_L.
\]
If $\tau=\{3^LX\}$ and $\tau_i=\{3^Lc_iX\}$, then exactly
\begin{equation}\label{real:eq:beta}
 \beta_{i,L}=c_i\tau-\tau_i
 \in q_i^{-1}\Z\cap(-1,p_i/q_i)=:\mathcal B_i.
\end{equation}
This is a fixed finite set of cardinality $p_i+q_i-1$. Every element
belongs to $\Z_3$ because $3\nmid q_i$.

Now $T_L-U_L$ is an integer multiple of $3^L$, so in $\Z_3$
\[
 c_iU_L+\beta_{i,L}
 =T_{i,L}-c_i(T_L-U_L)\equiv T_{i,L}\pmod{3^L}.
\]
Consequently
\begin{equation}\label{real:eq:membership}
 U_L\in W_L(\beta_{1,L},\ldots,\beta_{K,L}).
\end{equation}
The vector of shifts may change with $X$ and with $L$. We take the finite
union over $\prod_i\mathcal B_i$ at each depth; no consistency of a
single shift vector across depths is needed.

For a fixed integer part of $X$, the value $U_L$ determines a real
half-open cell of length $3^{-L}$. The first $L$ fractional digits form
the ordinary base-three integer $U_L$ in their usual written order;
the associated 3-adic word reads these same digits from low to high.
Both orders preserve the condition that every digit is $0$ or $1$.
There are at most $M+1$ integer parts. Thus the
number of occupied real cells of this length is at most
\begin{equation}\label{real:eq:real-count}
 (M+1)\left(\prod_{i=1}^K(p_i+q_i-1)\right)C_0\,3^{sL},
 \qquad L\ge L_0.
\end{equation}
Closed intervals give a cover with the same bound up to a fixed factor.
The factors preceding $3^{sL}$ are independent of $L$, so the upper
box dimension is at most $s$. The discrete scales $3^{-L}$ suffice,
since consecutive scales have bounded ratio.
\end{proof}

\begin{remark}\label{real:rem:changing}
The variability of the shifts is real, even for fully admissible
numbers. For $X=1/2$ and $c=1/4$, both $X$ and $cX=1/8$ belong to
$\CR$, while \eqref{real:eq:beta} equals $-1/4$ for odd $L$ and $0$ for
even $L$. Applying the container theorem after fixing each vector in
the finite union is what preserves prefix dependence. The proof does
not feed an unread real tail into a causal relation.
\end{remark}

\section{Countable compact covers and the dimension conclusion}
\begin{proof}[Proof of Theorem~\ref{real:thm:main}]
Fix $\delta>0$. Choose first $k$ sufficiently large, and then $r$ so
large that
\[
             \frac1k+\frac{2^{k+1}\log2}{rk\log3}<\delta.
\]
Leave a positive margin in this inequality. Choose $b$ with $3^b>r-1$,
and set $q=b+k$, $Q,C,D$ as in Section~\ref{sec:real:finite}. For integers $u\to\infty$
use the thresholds and number of templates
\[
 T_j=u^{2^{j-1}-1}\quad(2\le j\le r-1),\qquad
 T_{\rm s}=u^{2^{r-1}-1},\qquad K=u^{2^r}.
\]
Every ratio in \eqref{real:eq:theta} tends to zero. Thus
$\theta,\epsilon_F,\eta\to0$, and we can fix one finite $u$ with
$\theta\le1/2$ and $s\le\delta$ in \eqref{real:eq:s}.
All these parameters are now fixed.

For $\lambda\in\ER$ the true visit set
\[
                  S(\lambda)=\{n\ge0:2^n\lambda\in\CR\}
\]
is infinite. Apply Lemma~\ref{real:lem:extraction} to this set, obtaining
true indices $n_0,n_1,\ldots,n_K$ and finite template data. Then
$X=2^{n_0}\lambda$ and $c_iX=2^{n_i}\lambda$ all belong to $\CR$.
For some positive integer $M$, therefore,
\[
                 X\in K_{\R}(\boldsymbol c;M).
\]
By Proposition~\ref{real:prop:finite} and Lemma~\ref{real:lem:transfer}, this
compact set has upper box dimension at most $s\le\delta$.

Let $\mathcal D$ range over the finite data consisting of the visit
indices, the integer exponents and the ordinary integer coefficients
in Lemma~\ref{real:lem:extraction}. There are only countably many such data.
We have proved the cover
\begin{equation}\label{real:eq:cover}
 \ER\subseteq\bigcup_{\mathcal D}\ \bigcup_{M\ge1}
       2^{-n_0}K_{\R}(\boldsymbol c(\mathcal D);M).
\end{equation}
No 3-adic accumulation point or limiting coefficient indexes this union.
Each set is compact, and multiplication by $2^{-n_0}$ is a real
similarity, so its upper box dimension is at most $\delta$.

The span $R$ and the rational heights $p_i,q_i$ may be unbounded across
the cover. This causes no problem: in \eqref{real:eq:real-count} the depth
$L$ tends to infinity for one fixed finite system at a time. All of
those quantities affect only its multiplicative counting constant.
Countable stability of packing dimension and the usual inequalities
$\dimH\le\dimP\le\dimB$ on each compact component now imply
\[
                         \dimH\ER\le\dimP\ER\le\delta.
\]
Since $\delta>0$ was arbitrary, both dimensions are zero.
\end{proof}

The real argument uses infinitely many visits in its extraction step;
no uniform finite-visit conclusion for the real sets is asserted here.
The $3$-adic finite-visit conclusion of Theorem~\ref{thm:main} was proved
separately using the uniform finite extraction in
Theorem~\ref{thm:templates}. A zero-dimensional exceptional set need
not be countable or empty. In particular, Theorem~\ref{real:thm:main}
does not settle whether $1\in\ER$, the remaining pointwise question
associated with the powers of two.

The positive-dimensional truncated exceptional set does not contradict
the transfer: a condition only on $\lfloor2^n\lambda\rfloor$ gives no
restriction on arbitrarily long fractional prefixes. In that setting
the membership assertion \eqref{real:eq:membership} is unavailable.

\section{Exact checks for the real transfer}
The real proof consists of the arguments above, the reproduced finite
causal-container proof in Appendix~\ref{app:container}, and the
published unit-equation and sumset theorems~\cite{ESS,DeVos}.
The finite inputs originating in version~1.0.0 are explicitly identified. The accompanying
exact-arithmetic program checks the new finite-word bridge using
rational numbers, including terminating and periodic ternary expansions.
It verifies the bounds on the shifts, reduction with rational unit
denominators, and membership of the resulting finite words. The current
finite test set contains $1845$ rational inputs and $49$ slopes, with
$632835$ checked identities, including $12215$ instances with fully
admissible input and output expansions. Negative and nonintegral shifts
are included, as is the alternating-shift example in
Remark~\ref{real:rem:changing}.

These computations detect errors in the finite identities; they do not
prove the dimension theorem or replace the causal-container argument.
ChatGPT was used in developing and editing the manuscript, implementing
the exact checks, and conducting internal AI-assisted reviews from
several mathematical perspectives. Such reviews do not constitute
external peer review or proof-assistant verification. Authorship of this
manuscript is attributed solely to Timofei Iuzhakov.


% Author: Timofei Iuzhakov.
% Article section. Requires amsmath, amssymb, amsthm and the
% theorem, lemma, proposition, corollary, remark environments.
% Bibliography key: Lagarias. All labels in this section start with metric:.

\section{Digital rigidity and generic absence of three visits}
\label{metric:section}

Write
\[
 \Sigma=\left\{\sum_{j\geq0}e_j3^j:e_j\in\{0,1\}\right\}
 \subset\mathbb Z_3,\qquad G=1+3\mathbb Z_3,\qquad
 A=\Sigma\cap\mathbb Z_3^\times=1+3\Sigma.
\]
For a unit $b\in\mathbb Z_3^\times$, define
\[
 \begin{aligned}
 E_b^{(3)}&=\{\lambda\in\mathbb Z_3:
       b^n\lambda\in\Sigma\text{ for at least three distinct }n\geq0\},\\
 E_b&=\{\lambda\in\mathbb Z_3:
       b^n\lambda\in\Sigma\text{ for infinitely many }n\geq0\}.
 \end{aligned}
\]
All Haar measures on compact groups are normalized to have mass one.

\begin{theorem}\label{metric:main}
For Haar-almost every $b\in\mathbb Z_3^\times$,
\[
                         E_b^{(3)}=\{0\}.
\]
In particular, $E_b=\{0\}$ for Haar-almost every unit $b$.
The exceptional null set is independent of both the point $\lambda$ and
the three visit times.
\end{theorem}

This is a varying-base statement. The fixed-base exceptional sets were
studied by Lagarias \cite{Lagarias}; in particular, his
Theorem~1.6(iii) (Theorem~1.5(3) in the arXiv version~4) gives
$\dim_H E_2^{(3)}\geq (\log_3 2)/6$.
Thus the conclusion of Theorem~\ref{metric:main} cannot be imposed on the
prescribed base $2$. We do not attribute a varying-base conjecture to
Lagarias. Our proof uses a finite classification of extremal digital
planes, followed by an exact averaging identity over complete collections
of prefixes. Both ingredients are proved below.

\subsection{Full digital planes}

For $m\geq0$, put
\[
 D_m=\left\{\sum_{j<m}e_j3^j:e_j\in\{0,1\}\right\},\qquad
 B_m=\left\{\sum_{j<m}\sigma_j3^j:\sigma_j\in\{-1,1\}\right\},
 \qquad R_m=\frac{3^m-1}{2}.
\]
For unit coefficients $\alpha_1,\alpha_2,\alpha_3$ modulo $3^m$, let
\[
 P_m(C;\alpha)=\#\{(X_1,X_2,X_3)\in B_m^3:
               \alpha_1X_1+\alpha_2X_2+\alpha_3X_3\equiv C\pmod{3^m}\}.
\]
At most three of the eight next-digit triples satisfy a prescribed
residue modulo $3$. Consequently $P_m(C;\alpha)\leq3^m$.
We call a plane \emph{full} if equality holds. Equivalently, every
admissible node above the last level has three children.
Coordinate sign changes, permutations, and simultaneous multiplication
of the coefficients and target by a unit preserve the count.

\begin{theorem}[Finite digital rigidity]\label{metric:classification}
For $m\geq1$, a plane is full if and only if, after the operations just
listed, its coefficients and target have one of the following forms
modulo $3^m$:
\begin{align*}
 &(1,1,1),&& C\in B_m;\tag{a}\label{metric:classa}\\
 &(1,1,1+\tau3^g),&& C\in B_g,
       \quad 1\leq g<m,\quad\tau\in\{-1,1\}.\tag{b}\label{metric:classb}
\end{align*}
In (b), $B_g$ means its $2^g$ ordinary integer representatives reduced
modulo $3^m$, not all lifts of $B_g$ from modulus $3^g$.
\end{theorem}

We first establish the coefficient coincidence that underlies the
classification.

\begin{lemma}\label{metric:coincidence}
Normalize the coefficients by signs so that all three are $1$ modulo
$3$. If a depth-$m$ plane is full, two of these normalized coefficients
are equal modulo $3^m$.
\end{lemma}

\begin{proof}
Divide by the first coefficient. Use binary coordinates through the
bijection $X=2x-R_m$ from $D_m$ to $B_m$; modulo $3^m$ this is
$X\equiv2x+1/2$. For fixed coefficients let $F_l$ be the full binary
targets at depth $l$. Reduction sends $F_l$ into $F_{l-1}$, and each
of the three admitted first-digit choices below a full target has a
child target in $F_{l-1}$.

If all coefficients equal $1$ modulo $3^m$, there is nothing to prove.
Otherwise let $k<m$ be their first unequal depth, so that
\[
 \alpha_i=1+3^k a_i,\qquad a_1=0,
 \qquad a_1,a_2,a_3\text{ are not all equal modulo }3.
\]
Here $k\geq1$. For three equal binary coefficients, a target has three
first children precisely when its first digit is $1$ or $2$; all three
children have the same residual target. Induction gives
\[
                         F_k=R_k+D_k.
\]
For the canonical representative $V_0\in F_k$, every admitted triple
of depth-$k$ prefixes satisfies $x_1+x_2+x_3=V_0$ as an ordinary
equality: the sum is at most $3R_k<V_0+3^k$, and a negative wraparound
is impossible since $0\leq V_0<3^k$.

Consider a lift $V_0+3^k v$. Its next required residue, for an old
prefix $(x_1,x_2,x_3)$, is
$v-\sum_i a_i(x_i\bmod3)$. This residue must be nonzero.
If $V_0\equiv1\pmod3$, the three first-digit choices exclude
$v\in\{a_1,a_2,a_3\}$; if $V_0\equiv2\pmod3$, they exclude
$v\in\{a_1+a_2+a_3-a_j:1\leq j\leq3\}$.
All three choices occur in the full prefix tree. Three distinct
$a_i$ residues therefore permit no full lift. A split into two equal
residues and one other permits exactly one lift. Since a full
depth-$m$ target exists, the latter case applies, and $F_{k+1}$ is
injective under reduction modulo $3^k$. Fix the equal pair $i,j$;
then $\alpha_i\equiv\alpha_j\pmod{3^{k+1}}$.

This target injectivity persists at every depth $l\geq k+1$.
Indeed, suppose it holds at $l-1$, and take $V,V'\in F_l$ with
the same reduction modulo $3^k$. Their reductions modulo
$3^{l-1}$ agree. Choose the same admitted first-digit triple for
both targets. The child targets differ by $(V-V')/3$, which is zero
modulo $3^{l-2}$, and hence modulo $3^k$ when $l\geq k+2$.
Injectivity in $F_{l-1}$ gives $V\equiv V'\pmod{3^l}$.

If $m=k+1$, the selected coefficients already coincide modulo $3^m$.
Otherwise, among the three first children of a full depth-$m$ target
there are two whose targets differ by
$\pm(\alpha_i-\alpha_j)/3$. They agree modulo $3^k$.
Injectivity in $F_{m-1}$ forces their equality modulo $3^{m-1}$,
and therefore $\alpha_i\equiv\alpha_j\pmod{3^m}$.
\end{proof}

\begin{proof}[Proof of Theorem~\ref{metric:classification}]
By Lemma~\ref{metric:coincidence}, we may work with balanced
coefficients $(1,1,t)$, where $t\equiv1\pmod3$.
If $t\equiv1\pmod{3^m}$, every admitted child map is
$C\mapsto(C-\sigma)/3$, where $\sigma\in\{-1,1\}$ represents
$C\bmod3$. The plane is full exactly when $C\in B_m$.

Otherwise set $g=v_3(t-1)<m$, using any unit lift of $t$ to
$\mathbb Z_3$. The child maps of a full state are
\begin{equation}\label{metric:children}
 C\longmapsto\frac{C-\sigma t}{3}\quad\text{(twice)},
 \qquad
 C\longmapsto\frac{C-\sigma(2-t)}{3}\quad\text{(once)}.
\end{equation}
Put
\[
 \delta=\frac{2(1-t)}t,\qquad
 d\in\{-1,1\},\quad d\equiv\delta/3^g\pmod3,
 \qquad\epsilon=-d,\qquad D=1-\epsilon3^g.
\]
For any full state of depth $l$, always following the first map in
\eqref{metric:children} shows that $z=C/t$ has balanced digits
$\sigma_0,\ldots,\sigma_{l-1}$. Change only the branch at time $j$
to the second map and then follow first maps again. The signs through
time $j$ are unchanged. Telescoping the digit expansions gives
\[
                     z-\delta\sigma_j3^j\in B_l\pmod{3^l}.
\]
When $j+g<l$, the first affected digit is at $j+g$. It cannot be zero
modulo $3$, so
\begin{equation}\label{metric:period}
                         \sigma_{j+g}=\epsilon\sigma_j.
\end{equation}
It follows, at every finite depth $l\geq g$, that
\begin{equation}\label{metric:periodicform}
                  z\equiv P/D\pmod{3^l}\quad\text{for some }P\in B_g.
\end{equation}
At $l=g$ this only records the first $g$ digits, since $D\equiv1$
modulo $3^g$; at greater depths it is precisely
\eqref{metric:period}. No infinite extension has been assumed.

Apply \eqref{metric:periodicform} to the two different children of
the depth-$m$ root. Their normalized targets satisfy
\[
 z_0\equiv P_0/D,\qquad z_1\equiv P_1/D\pmod{3^{m-1}},
 \qquad P_0,P_1\in B_g,
\]
and, writing $\sigma=\sigma_0$,
\begin{equation}\label{metric:childdiff}
               P_1-P_0\equiv-\sigma\delta D/3\pmod{3^{m-1}}.
\end{equation}
This difference has valuation $g-1$. Thus the two balanced words
have equal lower $g-1$ digits and opposite highest digits.
The highest digit of $P_0$ is the root digit $\sigma_g$, which
equals $\epsilon\sigma$ by \eqref{metric:period}. Therefore, as
an ordinary integer equality,
\[
                    P_1-P_0=-2\epsilon\sigma3^{g-1}.
\]
Equation~\eqref{metric:childdiff} yields
$\delta D\equiv2\epsilon3^g\pmod{3^m}$. The identity
\[
                     \delta D-2\epsilon3^g=2(D-t)/t
\]
then gives $t\equiv D\pmod{3^m}$, and
\eqref{metric:periodicform} at the root gives $C\in B_g$ at the
full modulus $3^m$. This also proves the boundary case $m=g+1$:
the two child words have length $g$, and their differing digit is
still at position $g-1$.

Conversely, take exact representatives $t=1-\epsilon3^g$ and
$C\in B_g$. The maps \eqref{metric:children} become
\[
           \frac{C-\sigma}{3}\mathbin{\pm}\epsilon\sigma3^{g-1}.
\]
Both remain in $B_g$: they shift its lower $g-1$ remaining digits
and insert an allowed highest sign. Every reachable state is a unit
and has three children. The plane is full at all depths, and reduction
gives (b). The common-coefficient case was already proved.
\end{proof}

For later use define the two-input count
\begin{equation}\label{metric:Q}
 \begin{aligned}
 Q_m(U;a,b)&=\#\{(x,y)\in D_m^2:
                     U+ax+by\bmod3^m\in D_m\},\\
 H_m(a,b)&=\max_{U\bmod3^m}Q_m(U;a,b),
 \end{aligned}
\end{equation}
where $a,b$ are units. The balanced-coordinate conversion gives
\begin{equation}\label{metric:conversion}
 Q_m(U;a,b)=P_m\bigl((a+b-1)/2-2U;(a,b,-1)\bigr).
\end{equation}
In particular, $H_m(a,b)\leq3^m$.

\begin{corollary}\label{metric:inverse}
If $a,b\in\mathbb Z_3^\times$ and $H_m(a,b)=3^m$ for every
$m\geq1$, then
\[
                   a\in\{1,-1\},\quad\text{or}\quad
                   b\in\{1,-1\},\quad\text{or}\quad a\in\{b,-b\}.
\]
\end{corollary}
\begin{proof}
At every depth Lemma~\ref{metric:coincidence} gives one of these six
equalities modulo $3^m$. One equality occurs at unbounded depths,
and hence holds in $\mathbb Z_3$.
\end{proof}

\begin{corollary}[Enumeration and a rational height bound]
\label{metric:effective}
There are exactly $24m-20$ ordered unit pairs $(a,b)$ modulo $3^m$
for which $H_m(a,b)=3^m$, and exactly $28\cdot2^m-48$ full triples
$(a,b,U)$.
Moreover, let $a=A/D$, $b=B/D$, where $A,B,D$ are ordinary integers
prime to $3$ and $\max(|A|,|B|,|D|)\leq H$. If the pair is not full
at every depth, then
\[
                   3^m>4H^2\quad\Longrightarrow\quad
                   H_m(a,b)\leq3^m-1.
\]
\end{corollary}
\begin{proof}
Normalize $a,b$ separately by signs to be $1$ modulo $3$.
Theorem~\ref{metric:classification} and \eqref{metric:conversion}
give the normalized pairs
\[
 (1,1),\qquad (t,1),\ (1,t),\ (t^{-1},t^{-1}),
 \qquad t=1\pm3^g,\quad1\leq g<m.
\]
They are distinct: each nonconstant triple has a unique coincident
pair, and $v_3(t-1)=g$ separates the levels. There are $6m-5$
normalized pairs and four sign choices. The common pair has $2^m$
targets per sign choice; each nonconstant pair of level $g$ has
$2^g$. Hence the target count is
$4\cdot2^m+24\sum_{g=1}^{m-1}2^g=28\cdot2^m-48$.

Suppose now $3^m>4H^2$ and the rational pair is full at depth $m$.
Finite coefficient coincidence makes one of $A\pm B$, $A\pm D$,
$B\pm D$ divisible by $3^m$. Its absolute value is at most $2H$,
so it is zero. After normalization the triple is $(1,1,T/S)$
with $|S|,|T|\leq H$. If $T=S$, the game is full at all depths.
Otherwise put $g=v_3(T-S)$; then $3^g\leq2H<3^m$.
Classification gives $T-S(1\pm3^g)\equiv0\pmod{3^m}$, while
\[
                 |T-S(1\pm3^g)|\leq2H+2H^2\leq4H^2<3^m.
\]
This is again an exact equality, placing the coefficients in an
all-depth full family. The assumed nonfullness proves the implication.
\end{proof}

\subsection{Complete prefix cubes and Haar measure}

We use the elementary power operation $\rho^N$, for $\rho\in G$ and
$N\in\mathbb Z_3$. It is defined by the convergent binomial series,
or equivalently by $\exp(N\log\rho)$. The series for logarithm and
exponential give inverse group isomorphisms between $G$ and
$3\mathbb Z_3$. In particular a unit power is an automorphism of $G$.
The coefficients $\binom Nk$ are integral: approximate $N$ by
nonnegative integers in $\mathbb Z_3$ and use continuity of each
binomial polynomial. These series also show that
$\rho\mapsto\rho^N$ is $1$-Lipschitz. We will use these facts only
on principal units.

For fixed $N\in\mathbb Z_3\setminus\{0,1\}$ define
\[
 \mathcal B_N=\{\rho\in G:\text{some }\mu\in A\text{ satisfies }
                           \rho\mu,\rho^N\mu\in A\}.
\]
This set is compact, being the projection of a closed subset of
$G\times A$.

\begin{proposition}\label{metric:fixedratio}
For every fixed $N\in\mathbb Z_3\setminus\{0,1\}$, the set
$\mathcal B_N$ has Haar measure zero in $G$.
\end{proposition}

\begin{proof}
If $N\equiv2\pmod3$, then $\mathcal B_N\subseteq\{1\}$.
Indeed, write $\rho=1+3^g t\ne1$ with $g\geq1$ and $t$ a unit.
The three numbers $\mu,\rho\mu,\rho^N\mu$ have equal digits below
$g$, whereas their digits at $g$ differ by $0,t,2t$ modulo $3$
because $\mu\equiv1\pmod3$. One of these three digits is $2$.

Next suppose $N=3^s r$ with $s\geq1$ and $r$ a unit. Set
\begin{equation}\label{metric:coefficients}
 a(\rho)=(1-N)\rho^N,\qquad b(\rho)=r\rho^{N-1},\qquad
 \Phi(\mu,\nu)=\mu^{1-N}\nu^N.
\end{equation}
Both coefficients are units. A level-$q$ node is a pair
$(\mu,\nu)\in A_q^2$, where $A_q=1+3D_{q-1}$, for which
$\Phi(\mu,\nu)\bmod3^q\in A_q$. Representatives in $A_q$ are
ordinary nonnegative integers. Its ratio is $\rho=\nu/\mu\in G$.

Fix $m\geq1$ and $q\geq m+s+1$. Group pairs by fixing $\mu$ and
the digits of $\nu$ below position $q-s$. Write the resulting pairs as
\[
            (\mu,\nu_e),\qquad
            \nu_e=\nu_0+3^{q-s}e,\qquad e\in D_s.
\]
Every group containing a node contains all its $2^s$ possible nodes.
More precisely, with $\rho_0=\nu_0/\mu$,
\begin{equation}\label{metric:cubeshift}
 \Phi(\mu,\nu_e)\equiv\Phi(\mu,\nu_0)
               +3^q b(\rho_0)e\pmod{3^{q+m+1}}.
\end{equation}
To prove this, expand in $3^{q-s}e/\nu_0$. From
$\binom Nk=(N/k)\binom{N-1}{k-1}$ we obtain
$v_3\binom Nk\geq s-v_3(k)$. For $k\geq2$ the valuation of the
$k$th term is at least
\[
 k(q-s)+s-v_3(k)\geq2q-s\geq q+m+1;
\]
the first inequality follows from $v_3(k)\leq k-2$ and $q-s\geq1$.
This proves \eqref{metric:cubeshift}. In addition, all ratios in the
group agree modulo $3^{m+1}$.

Let $U$ be the block of $m$ digits of $\Phi(\mu,\nu_0)$ beginning
at position $q$. The descendants of a group have the form
\[
                     (\mu+3^q x,\nu_e+3^q y),\qquad x,y\in D_m.
\]
Integral binomial coefficients and $2q\geq q+m+1$ give their exact
new-block condition
\begin{equation}\label{metric:cubelinear}
              U+be+ax+3^sby\bmod3^m\in D_m,
\end{equation}
where $a,b$ are evaluated at $\rho_0$. They can be used throughout
the group by the Lipschitz property and the common ratio residue.
The lower $q$ digits remain admissible.

As $e\in D_s$ and $y\in D_m$ vary, $z=e+3^s y$ runs bijectively
over $D_{m+s}$. Each residue in $D_m$ has exactly $2^s$ such lifts
modulo $3^m$, also when $m<s$. Summing \eqref{metric:cubelinear}
therefore gives the exact descendant count
\begin{equation}\label{metric:cubecount}
                \#\{\text{descendants of the group}\}
                         =2^s Q_m(U;a,b).
\end{equation}

Fix a ratio class $\mathcal C\subset G$ modulo $3^{m+1}$.
The coefficients are constant modulo $3^m$ on this class. Write
$H_{\mathcal C}=H_m(a,b)$, and let $T_q(\mathcal C)$ count its
level-$q$ nodes. All the preceding groups stay in $\mathcal C$,
as do their descendants. Since each group has $2^s$ original nodes,
\eqref{metric:cubecount} implies
\begin{equation}\label{metric:contraction}
             T_{q+m}(\mathcal C)\leq H_{\mathcal C}T_q(\mathcal C),
             \qquad q\geq m+s+1.
\end{equation}
If $H_{\mathcal C}<3^m$, iterate with $q_j=q_0+jm$ to obtain
$T_{q_j}(\mathcal C)\leq T_{q_0}(\mathcal C)H_{\mathcal C}^j$.

For a ratio $\rho$ modulo $3^q$, let $n_q(\rho)$ count the
$\mu\in A_q$ for which $\rho\mu$ and $\rho^N\mu$ lie in $A_q$
modulo $3^q$. The substitution $\nu\equiv\rho\mu\pmod{3^q}$
is a bijection, and $G$ has $3^{q-1}$ residue classes modulo $3^q$.
Consequently
\[
 \int_{\mathcal C}n_q(\rho)\,d\rho
                     =\frac{T_q(\mathcal C)}{3^{q-1}}.
\]
Every $\rho\in\mathcal B_N$ has $n_q(\rho)\geq1$ at each depth.
It follows that
\[
 \operatorname{Haar}(\mathcal B_N\cap\mathcal C)
 \leq\frac{T_{q_j}(\mathcal C)}{3^{q_j-1}}\longrightarrow0
 \qquad(H_{\mathcal C}<3^m).
\]
Taking the countable union over finite windows and ratio classes
reduces the proof to the residual set
\[
 \mathcal F_N=\{\rho\in G:H_m(a(\rho),b(\rho))=3^m
                                      \text{ for every }m\geq1\}.
\]

In fact $\mathcal F_N$ contains at most two points. Let
$\sigma\in\{-1,1\}$ represent $r\bmod3$. Since
$\rho^N\equiv1\pmod{3^{s+1}}$,
\[
                    a(\rho)-1\equiv-N\pmod{3^{s+1}},
                    \qquad v_3(a(\rho)-1)=s.
\]
Thus $a\ne1$, and $a\equiv1\pmod3$ also excludes $a=-1$.
Corollary~\ref{metric:inverse} leaves only $b=\sigma$ or
$a=\sigma b$. The first has the unique solution
$\rho_1=(\sigma/r)^{1/(N-1)}\in G$, because $N-1$ is a unit.
The second, using $a/b=((1-N)/r)\rho$, has the unique solution
$\rho_2=\sigma r/(1-N)\in G$. Hence
$\mathcal F_N\subseteq\{\rho_1,\rho_2\}$, proving nullity.
No assertion that either candidate is full is needed.

Finally, if $N\equiv1\pmod3$, exchange $\mu$ and $\nu=\rho\mu$.
The ratio becomes $\rho^{-1}$ and the exponent becomes $1-N$,
which is divisible by $3$ and different from $0,1$.
Thus inversion maps $\mathcal B_N$ into $\mathcal B_{1-N}$.
Inversion preserves Haar measure, and the preceding case applies.
\end{proof}

\subsection{One null set for all points and times}

\begin{proof}[Proof of Theorem~\ref{metric:main}]
For every nonzero integer $e$, the map $P_e:G\to G$ given by
$P_e(\beta)=\beta^e$ has null preimages of Haar-null sets.
Indeed, in logarithmic coordinates it is $x\mapsto ex$ on
$3\mathbb Z_3$, an injective homomorphism onto an open subgroup.
Its pushforward measure is normalized Haar measure on that subgroup,
and is therefore absolutely continuous with respect to Haar measure
on $G$.

By Proposition~\ref{metric:fixedratio}, the set
\[
 \mathcal N=
 \bigcup_{e\in\mathbb Z\setminus\{0\}}
 \ \bigcup_{N\in(\mathbb Q\cap\mathbb Z_3)\setminus\{0,1\}}
                         P_e^{-1}(\mathcal B_N)
\]
is Haar-null. Fix $\beta\in G\setminus\mathcal N$. If
$\lambda\ne0$ had three distinct visit times $n_0,n_1,n_2$, put
$v=v_3(\lambda)$. Interchange $n_1,n_2$ if necessary so that
\[
 v_3(e)\leq v_3(e'),\qquad e=n_1-n_0,\qquad e'=n_2-n_0.
\]
Then $N=e'/e\in(\mathbb Q\cap\mathbb Z_3)\setminus\{0,1\}$,
and the quantities
\[
             \mu=3^{-v}\beta^{n_0}\lambda,\qquad\rho=\beta^e
\]
satisfy $\mu,\rho\mu,\rho^N\mu\in A$.
The identity $\rho^N=\beta^{e'}$ follows from logarithms.
Thus $\rho\in\mathcal B_N$, contradicting the choice of $\beta$.
This proves the theorem on $G$ simultaneously for all nonzero
$\lambda$ and all triples of times.

For $b\equiv-1\pmod3$, the visit times of a fixed nonzero point
have a common parity: the leading unit digit of $b^n\lambda$ must
be $1$. Three visits therefore give three visits of $b^{r_0}\lambda$
under the base $\beta=b^2\in G$, for some $r_0\in\{0,1\}$.
Squaring maps the coset $-1+3\mathbb Z_3$ bijectively onto $G$
and has null preimages of null sets; write $b=-u$ and use the unit
power automorphism $u\mapsto u^2$. This proves the statement on
all units. Finally, $0$ always visits $\Sigma$ and
$E_b\subseteq E_b^{(3)}$.
\end{proof}

\begin{remark}\label{metric:scope}
The averaging in \eqref{metric:cubecount} is over a complete group
of prefixes. It does not give a contraction on a prescribed ratio
fiber. Accordingly, neither Theorem~\ref{metric:main} nor the
effective finite-window bound in Corollary~\ref{metric:effective}
asserts finiteness of visits for a specified base and a specified
point. The proofs above are independent of finite enumeration;
the accompanying exact checks verify small instances of the
classification and the complete-cube identity.
\end{remark}

% Author: Timofei Iuzhakov
% Integration: amsmath, amssymb, amsthm, mathrsfs; theorem, lemma, proposition,
% remark and proof environments. Bibliography keys: Lagarias, KP2005,
% AL2014. No other local file is a mathematical dependency.
\section{Summable denominators and positive-density multipliers}
\label{sec:positive-density}

Write
\[
 \mathcal D=\left\{\sum_{j=0}^{L-1}\epsilon_j3^j:
       L\geq1,\ \epsilon_j\in\{0,1\}\right\},\qquad
 \Sigma=\left\{\sum_{j\geq0}\epsilon_j3^j:
       \epsilon_j\in\{0,1\}\right\}\subset\mathbb Z_3.
\]
The first sum is finite and defines an ordinary nonnegative integer;
the second converges in the $3$-adic metric. For a positive integer $M$ set
\[
 \begin{aligned}
 C_M&=\{x\in\Sigma:Mx\in\Sigma\},\\
 \mathcal M_C&=\{M\geq1:\exists a\in\mathcal D\setminus\{0\},\ Ma\in\mathcal D\},\\
 \mathcal M_H&=\{M\geq1:\dim_H C_M>0\}.
 \end{aligned}
\]
Here Hausdorff dimension is taken in the $3$-adic metric, and
$M_C(X)=\#(\mathcal M_C\cap[1,X])$, with analogous notation for $M_H$.
These are the two multiplier-counting questions in Problems~1 and~2 of
the arXiv version of \cite{Lagarias} (journal Problems~1.10 and~1.11).
The nonzero requirement in $\mathcal M_C$ is essential.

\begin{theorem}\label{thm:pd-main}
There are constants $c_1,c_2>0$ such that, for all sufficiently large $X$,
\[
 M_H(X)\geq c_1X,\qquad M_H(X)-M_C(X)\geq c_2X.
\]
In particular $M_H(X)\asymp X$. This theorem does not assert that
$M_H(X)/X$ has a limit, nor that this limit is $1/2$.
\end{theorem}

The proof uses reduced denominators in the ordinary middle-thirds set
\[
 \mathcal C=\left\{\sum_{j\geq1}2\epsilon_j3^{-j}:\epsilon_j\in\{0,1\}\right\},
 \qquad s=\frac{\log2}{\log3}.
\]
Define its set of internal reduced denominators prime to $3$ by
\[
 \mathscr D=\{d>1:3\nmid d,\ \exists a,\ 0<a<d,\ (a,d)=1,
                                      \ a/d\in\mathcal C\}.
\]
Thus endpoints and unreduced representations are not counted.

\begin{theorem}\label{thm:pd-denominators}
Put $\kappa=4/3-2s>0$. As $Y\longrightarrow\infty$,
\[
 \sum_{\substack{d\in\mathscr D\\d>Y}}\frac1d
       \ll (\log Y)^{-9}+Y^{-\kappa}.
\]
In particular $\sum_{d\in\mathscr D}1/d$ converges.
\end{theorem}

\subsection{The carry obstruction}

We give the carry argument explicitly; finite carry presentations of
these intersections also appear in \cite{AL2014}.
Suppose $M>1$ and $M\equiv1\pmod3$. Reading input and output digits
$a,b\in\{0,1\}$ from low to high gives the transition
\begin{equation}\label{eq:pd-carry}
 3c'=c+Ma-b,
\end{equation}
starting at $c=0$. The allowed pairs $(a,b)$, according as
$c\equiv0,1,2\pmod3$, are respectively
\[
 \{(0,0),(1,1)\},\qquad\{(0,1)\},\qquad\{(1,0)\}.
\]
All reachable carries are nonnegative. After $t$ steps, if $A,B$ are
the input and output integers read so far, then $3^tc=MA-B$;
since $A\leq(3^t-1)/2$ and $B\geq0$, every reachable carry satisfies
$c<M/2$. The reachable graph is therefore finite and has no dead ends.
Conversely, an infinite path gives $Mx=y$ in $\mathbb Z_3$, because its
bounded carry error is divisible by $3^t$ at every stage.

\begin{lemma}\label{lem:pd-terminal}
If $M>1$, $M\equiv1\pmod3$, and $\dim_H C_M=0$, then there is an integer
$c$ with $0\leq2c<M$ such that
\[
 \frac{2c+1}{M+1}\in\mathcal C\cap(0,1).
\]
\end{lemma}
\begin{proof}
If a reachable strongly connected component contains a vertex with two
outgoing edges inside that component, choose closed paths $U,V$ based
there and starting with these different edges. Their first input digits
differ. The return words $UV,VU$ have the same length and different
first input digits. Arbitrary concatenations of them, after a fixed
path from the initial state, yield a binary block Cantor set of positive
Hausdorff dimension. Thus no such component exists when $\dim_H C_M=0$.

Choose a terminal strongly connected component. All its outgoing edges
remain inside it, so its vertices cannot be $0$ modulo $3$. Each vertex
has exactly one outgoing edge, and the component is a directed cycle.
On this cycle $a+b=1$. Put $q=M+1$ and $r_j=(2c_j+1)/q$. Equation
\eqref{eq:pd-carry} becomes $3r_{j+1}=r_j+2a_j$. If the cycle length is
$L$, iteration gives
\[
 r_0=\frac{2\sum_{j=0}^{L-1}a_j3^j}{3^L-1}\in\mathcal C.
\]
The inequalities on the carries give $0<r_0<1$.
\end{proof}

Only this necessary condition for zero dimension will be used. A
terminal cycle can coexist with a different branching component.

\begin{lemma}\label{lem:pd-inclusion}
$\mathcal M_C\subseteq\mathcal M_H$.
\end{lemma}
\begin{proof}
Let $Ma=b$ with positive $a,b\in\mathcal D$, and choose $L$ at least
the ternary lengths of both. All points
$x=a\sum_{j\geq0}\epsilon_j3^{Lj}$, $\epsilon_j\in\{0,1\}$, lie in
$C_M$: the input blocks are $0,a$ and the output blocks are $0,b$,
without overlap. The set of the sums has dimension
$\log2/(L\log3)>0$; multiplication by the nonzero integer $a$ is a
$3$-adic similarity and preserves dimension.
\end{proof}

\subsection{A small-order estimate}

For $n\geq2$ with $(n,3)=1$ let $\operatorname{ord}_n(3)$ be the multiplicative order
of $3$ modulo $n$, and let $\lambda(n)$ be Carmichael's function.
The only non-elementary input in this section is the following
unconditional estimate, stated as Lemma~10 of \cite{KP2005}:
\begin{equation}\label{eq:pd-kp}
 \#\{n\leq X:\lambda(n)\leq n\exp(- (\log\log n)^3)\}
       \ll\frac{X}{(\log X)^{10}}.
\end{equation}
Finitely many small $n$ can be omitted in this display. No hypothesis
such as the generalized Riemann hypothesis is assumed.

\begin{lemma}\label{lem:pd-small-order}
\[
 \#\{2\leq n\leq X:(n,3)=1,\ \operatorname{ord}_n(3)<n^{1/3}\}
       \ll X(\log X)^{-10}.
\]
\end{lemma}
\begin{proof}
Let $\mathscr P$ be the primes $p\ne3$ for which
$\operatorname{ord}_p(3)\leq p^{2/5}$. Every such $p\leq z$ divides
$\prod_{1\leq j\leq z^{2/5}}(3^j-1)$. Taking logarithms gives
$\#\{p\in\mathscr P:p\leq z\}=O(z^{4/5})$, and partial summation
gives $\sum_{p\in\mathscr P}p^{-9/10}<\infty$.

Let $n_{\mathscr P}$ be the largest divisor of $n$ supported on
$\mathscr P$. For $f(n)=n_{\mathscr P}^{1/10}$, the multiplicative
function $f*\mu$ is nonnegative. Hence
\[
 \sum_{n\leq X} f(n)
 \leq X\sum_{d\geq1}\frac{(f*\mu)(d)}d
 =X\prod_{p\in\mathscr P}\frac{1-p^{-1}}{1-p^{-9/10}}
 \ll X.
\]
It follows that $n_{\mathscr P}>X^{1/20}$ for at most $O(X^{199/200})$
integers $n\leq X$.

Similarly put $R(n)=n/\operatorname{rad}(n)$ and $h(n)=R(n)^{1/2}$.
Here $(h*\mu)(p)=0$, $(h*\mu)(p^e)\geq0$, and
$\sum_{e\geq2}(h*\mu)(p^e)/p^e=O(p^{-3/2})$.
The same argument gives $\sum_{n\leq X}h(n)\ll X$, so
$R(n)>X^{1/20}$ for at most $O(X^{39/40})$ integers $n\leq X$.

We also need the elementary inequality
\begin{equation}\label{eq:pd-order-product}
 \operatorname{ord}_n(3)\geq
       \frac{\lambda(n)}n\prod_{p\mid n}\operatorname{ord}_p(3)
       \qquad ((n,3)=1).
\end{equation}
Indeed, for each $p^e\Vert n$ write
$\lambda(p^e)=a_pb_p$ with $a_p=\operatorname{ord}_p(3)$.
The integer $\operatorname{lcm}_p(a_pb_p)$ divides
$(\prod_p b_p)\operatorname{lcm}_p a_p$, while
$\operatorname{lcm}_p a_p$ divides $\operatorname{ord}_n(3)$ and
$\prod_p\lambda(p^e)\leq n$. Rearrangement proves
\eqref{eq:pd-order-product}.

Consider $X/2<n\leq X$, excluding the two moment exceptions above
and the exceptions in \eqref{eq:pd-kp}. The product of the distinct
prime factors of $n$ outside $\mathscr P$ is at least
\[
 \frac{n}{n_{\mathscr P}R(n)}\geq\frac12X^{9/10}.
\]
Consequently \eqref{eq:pd-order-product} gives
\[
 \operatorname{ord}_n(3)
 \geq\exp(- (\log\log X)^3)
                       \left(\tfrac12X^{9/10}\right)^{2/5}
 >X^{1/3}
\]
for sufficiently large $X$. In the last step one may use
$\exp(- (\log\log X)^3)\geq X^{-1/100}$ eventually, since
$9/25-1/100=7/20>1/3$. This proves the required estimate on each
dyadic interval. Sum over the intervals between $\sqrt X$ and $X$,
where their logarithms are comparable to $\log X$, and absorb the
at most $\sqrt X$ remaining integers.
\end{proof}

\subsection{Summability of Cantor denominators}

\begin{proof}[Proof of Theorem~\ref{thm:pd-denominators}]
Distinct reduced fractions of denominator at most $T$ are separated by
at least $T^{-2}$. Cover $\mathcal C$ by its $2^k$ level-$k$ intervals,
where $k=\lfloor2\log_3T\rfloor+1$. Each interval has length
$3^{-k}<T^{-2}$ and contains at most one of these fractions. Thus the
number of reduced rational points of $\mathcal C$ with denominator at
most $T$ is $O(T^{2s})$.

For each $d\in\mathscr D$, the orbit of one reduced $a/d\in\mathcal C$
under $x\mapsto3x\pmod1$ stays in $\mathcal C$ and contains exactly
$\operatorname{ord}_d(3)$ points. Orbits belonging to different
reduced denominators are disjoint. It follows that
\[
 \#\{d\in\mathscr D:T/2<d\leq T,
                       \ \operatorname{ord}_d(3)\geq d^{1/3}\}
       \ll T^{2s-1/3}.
\]
Lemma~\ref{lem:pd-small-order} handles the remaining denominators,
so
\[
 \#(\mathscr D\cap(T/2,T])
       \ll T(\log T)^{-10}+T^{2s-1/3}.
\]
Since $s<2/3$ (equivalently $8<9$), we have
$\kappa=1-(2s-1/3)>0$. Sum the reciprocal-denominator bound on
the shells $(2^jY,2^{j+1}Y]$, $j\geq0$. The power part is
$O(Y^{-\kappa})$ and the logarithmic part is
$O(\sum_{j\geq0}(\log Y+j)^{-10})=O((\log Y)^{-9})$.
\end{proof}

\subsection{A positive-density sifted set}

Let $\mathscr B$ be the odd members of $\mathscr D$ and put
\[
 \mathcal Q=\{q\geq1:q\equiv5\pmod6,
                   \ d\nmid q\text{ for every }d\in\mathscr B\}.
\]

\begin{lemma}\label{lem:pd-sieve}
The set $\mathcal Q$ has a natural density $\delta>0$, with
\[
 \delta\geq\frac16\prod_{d\in\mathscr B}(1-1/d)>0.
\]
\end{lemma}
\begin{proof}
For a finite subset $\mathscr B_0\subset\mathscr B$, its avoidance
set is periodic. All its moduli are coprime to $6$. By the Chinese
remainder theorem, conditioning on $q\equiv5\pmod6$ leaves uniform,
independent prime-power residue coordinates modulo
$\operatorname{lcm}(\mathscr B_0)$. The truncated prime valuations
are therefore independent. Each event $d\nmid q$ is decreasing in
these valuations. Such decreasing events are positively associated.
For completeness, in one coordinate, increasing or decreasing
functions $f,g$ in the same direction satisfy
\[
 2\operatorname{Cov}(f(U),g(U))
 =\mathbb E[(f(U)-f(U'))(g(U)-g(U'))]\geq0,
\]
where $U'$ is an independent copy. Conditioning on one coordinate
and induction prove the same inequality in a product distribution.
Applying it successively to products of decreasing indicators gives
a density at least $\frac16\prod_{d\in\mathscr B_0}(1-1/d)$.
This is a lower bound; divisibility by different composite moduli
is not being treated as independent.

Take $\mathscr B_0=\mathscr B\cap[1,Y]$, and denote the density of
this finite avoidance set by $\delta_Y$. The densities decrease to
some $\delta>0$, since Theorem~\ref{thm:pd-denominators} makes the
displayed infinite product positive. Uniformly in $X$, the count
removed on imposing all the conditions with $d>Y$ is at most
\[
 \sum_{\substack{d\in\mathscr B\\d>Y}}\lfloor X/d\rfloor
       \leq X\sum_{\substack{d\in\mathscr B\\d>Y}}1/d.
\]
First let $X\to\infty$ with $Y$ fixed, then let $Y\to\infty$.
The tail tends to zero, proving that $\mathcal Q$ itself has density
$\delta$ and establishing the claimed lower bound.
\end{proof}

\subsection{Positive dimension and separation}

\begin{proof}[Proof of Theorem~\ref{thm:pd-main}]
If $q\in\mathcal Q$, no internal fraction $a/q$ belongs to
$\mathcal C$: its reduced denominator would be an odd member $d>1$
of $\mathscr D$ dividing $q$. Set $M=q-1$. Then $M>1$ and
$M\equiv1\pmod3$, so Lemma~\ref{lem:pd-terminal}, by contraposition,
gives $M\in\mathcal M_H$. Hence $M_H(X)\geq\delta X+o(X)$.

If a positive Cantor integer has highest nonzero ternary position $j$,
its value belongs to $[3^j,(3^{j+1}-1)/2]$. Consequently
\[
 \{b/a:a,b\in\mathcal D\setminus\{0\}\}
   \subseteq\bigcup_{j\in\mathbb Z}
       \left(\tfrac23\,3^j,\tfrac32\,3^j\right).
\]
In particular no integer in $[3T/2,2T]$ belongs to $\mathcal M_C$
when $T$ is a power of $3$. Choose such a $T$ with $2T\leq X<6T$.
By Lemma~\ref{lem:pd-sieve}, the interval
$[3T/2+1,2T+1]$ contains $\delta T/2+o(T)$ elements of $\mathcal Q$.
Their translates $M=q-1$ are in
$(\mathcal M_H\setminus\mathcal M_C)\cap[1,X]$. For all sufficiently
large $X$ there are at least $\delta T/4>\delta X/24$ of them.
Use Lemma~\ref{lem:pd-inclusion} to identify this set difference
with the difference of the two counting functions. The trivial upper
bound $M_H(X)\leq\lfloor X\rfloor$ completes the assertion
$M_H(X)\asymp X$.
\end{proof}

\begin{remark}
The positive density proved above belongs to the explicit sifted
subfamily $\{q-1:q\in\mathcal Q\}$, not necessarily to all of
$\mathcal M_H$. The separation estimate is not a linear lower bound
for $M_C(X)$. These results concern varying integer multipliers; they
make no finiteness assertion about the visits of a fixed nonzero
$3$-adic point under multiplication by powers of $2$.
\end{remark}


\section*{New verification in version 3}
The new standard-library Python checks verify finite digital fibers,
all coefficient and target cases in the stated test ranges, complete
prefix-cube identities, and finite carry, rational-denominator and sieve
identities. These computations supplement the proofs for arbitrary
depths and parameters; they do not prove an asymptotic by extrapolation.
The exact new source files received fresh independent internal
AI-assisted reviews. The earlier proof sections and causal-container
appendix are retained from the released version 2 sources; this update
does not describe their preservation as a fresh line-by-line audit.
The metric branch uses no external unproved hypothesis. The counting
branch explicitly states its external analytic input,
Kurlberg--Pomerance Lemma 10~\cite{KP2005}; no GRH assumption is used.

\section*{Dependency and verification record}
For the dimension proofs retained from version 2, the imported nonstandard finite result is
Theorem~\ref{thm:container}, from Appendix A of the exact release
v1.0.0 of \cite{Core}, commit
\begin{center}\small\ttfamily
81c7f5ae45c3ce507173b82baca48fdfeb52c6ef
\end{center}
Its complete proof is included unchanged in Appendix~\ref{app:container}.
The SHA-256 digest of the appendix source
\texttt{causal\_container.tex} is
\begin{center}\small\ttfamily
1ae15961cd3bc705e7313c9e56de34e5127\\
7cffaa7cc3a58c6e4c9cb6e672750
\end{center}
with the two displayed lines concatenated. This identifies the exact
imported proof on which the finite entropy estimate depends. The
unit-equation theorem is the published result \cite{ESS}; finite
Ramsey theory is used only in its ordinary finite form. The real argument
also uses Kneser's addition theorem~\cite{DeVos}, through the full proof
of Lemma~\ref{vone:lem:capacity} reproduced from
\cite[Lemma~2.1]{Core}. The statements and proofs of the uniform finite
extraction and of the real infinite-set extraction are kept separate.
The present version does not assert a real finite-visit nesting theorem,
a finite-visit Assouad theorem, or an arbitrary-target generalization.

The arithmetic extraction and the entropy transfer were subjected to
independent internal AI-assisted reviews. Exact finite checks were
also used to test determinant identities and low-dimensional cases.
These checks do not replace the arguments for arbitrary $d$, all
Ramsey refinements, or the causal-container theorem.
AI assistance was used in developing proofs, writing the manuscript,
and implementing checks. This is not external peer review or
proof-assistant verification. Authorship is attributed solely to
Timofei Iuzhakov.

\begin{thebibliography}{99}
\bibitem{Lagarias}
J. C. Lagarias,
\emph{Ternary expansions of powers of 2},
J. London Math. Soc. (2) \textbf{79} (2009), 562--588.
\href{https://doi.org/10.1112/jlms/jdn080}{doi:10.1112/jlms/jdn080}.

\bibitem{ABL}
W. C. Abram, A. Bolshakov and J. C. Lagarias,
\emph{Intersections of multiplicative translates of 3-adic Cantor sets II:
two infinite families},
Experimental Math. \textbf{26} (2017), 468--489.
\href{https://doi.org/10.1080/10586458.2016.1205532}
{doi:10.1080/10586458.2016.1205532};
\href{https://arxiv.org/abs/1508.05967}{arXiv:1508.05967}.

\bibitem{ESS}
J.-H. Evertse, H. P. Schlickewei and W. M. Schmidt,
\emph{Linear equations in variables which lie in a multiplicative group},
Ann. of Math. (2) \textbf{155} (2002), 807--836.
\url{https://annals.math.princeton.edu/2002/155-3/p04}.

\bibitem{Core}
T. Iuzhakov,
\emph{Zero Hausdorff and packing dimension for the
$3$-adic exceptional set of powers of two},
research preprint, release v1.0.0 (5 November 2026), Lemma~2.1 and Appendix~A.
\url{https://github.com/ppadfoot/lagarias-zero-dimension/releases/tag/v1.0.0}.
\bibitem{LagariasArxiv}
J. C. Lagarias,
\emph{Ternary expansions of powers of 2}, arXiv manuscript.
\href{https://arxiv.org/abs/math/0512006}{arXiv:math/0512006}.

\bibitem{DeVos}
M. DeVos,
\emph{A short proof of Kneser's addition theorem for abelian groups}.
\url{https://www.sfu.ca/~mdevos/papers/kneser.pdf}.

\bibitem{AL2014}
W. C. Abram and J. C. Lagarias,
\emph{Intersections of multiplicative translates of $3$-adic Cantor sets},
J. Fractal Geom. \textbf{1} (2014), 349--390.
\href{https://arxiv.org/abs/1308.3133}{arXiv:1308.3133}.
\bibitem{KP2005}
P. Kurlberg and C. Pomerance,
\emph{On the period of the linear congruential and power generators},
Acta Arith. \textbf{119} (2005), 149--169.
\href{https://arxiv.org/abs/math/0405120}{arXiv:math/0405120v2}.
\bibitem{Quotients}
T. Iuzhakov, \emph{Integer quotients of ternary Cantor numbers:
a power lower bound and an infinite exceptional family}, research
preprint, version 1.0.0, 6 November 2026.
\href{https://github.com/ppadfoot/ternary-cantor-quotients/releases/tag/v1.0.0}{Versioned release}.
\end{thebibliography}
\appendix
\section{The finite causal-container theorem and its proof}
\label{app:container}
The following source is reproduced unchanged from Appendix A of
\cite{Core}, release v1.0.0. Its symbols are local to this appendix.

\subsection{Exact theorem}\label{exact-theorem}

Let \(r\ge3\), let \(V_0=\{0,\ldots,N-1\}\), let \(\Omega\) be a finite alphabet of size \(Q\), and let \(\mathcal W\subseteq\mathcal X^L\) be nonempty, where \(\mathcal X\) is any finite input alphabet. For every \(i\in V_0\), a word \(x\) determines a letter \(Z_i(x)\in\Omega\). No locality hypothesis on the maps \(Z_i\) is needed for this finite theorem.

Let \(\mathcal E\) be a fixed family of ordered positional edges

\[
e=(s=i_0<i_1<\cdots<i_{r-1}).
\]

For each edge there is a relation \(\mathcal R_e(x_{<s})\subseteq\Omega^r\), depending only on the indicated input prefix, and every \(x\in\mathcal W\) satisfies

\[
(Z_{i_0}(x),\ldots,Z_{i_{r-1}}(x))\in\mathcal R_e(x_{<s}).
\]

Assume:

\begin{enumerate}
\def\labelenumi{\arabic{enumi}.}
\tightlist
\item
  for a fixed source \(s\), the future position sets \(e\setminus\{s\}\) of distinct outgoing edges are disjoint;
\item
  every unordered pair of positions belongs to at most \(C\) positional edges.
\end{enumerate}

Fix positive integer thresholds \(T_{\rm s},T_2,\ldots,T_{r-1}\), and put

\[
D=CQ^{r-3},
\]

\[
\theta=Q\left[
\frac{D\sum_{j=2}^{r-1}T_j+1}{T_{\rm s}}
+\frac1{T_2}
+D\sum_{j=3}^{r-1}\frac{1+\sum_{l=2}^{j-1}T_l}{T_j}
\right].
\]

Assume \(\theta\le1/2\). Then each word can be assigned a set \(F(x)\) of at most \(\theta N\) pairs \((i,Z_i(x))\), with at most one pair for each position, and nonempty sets \(A_i(x)\subseteq\Omega\), such that:

\begin{itemize}
\tightlist
\item
  \(Z_i(x)\in A_i(x)\);
\item
  \(A_i(x)\) is determined by \((x_{<i},F(x))\);
\item
  except for at most
\end{itemize}

\[
Q\left(T_{\rm s}+D\sum_{j=2}^{r-1}T_j\right)N
\]

edges, the full product of the assigned alphabets is compatible:

\[
\prod_{\nu=0}^{r-1}A_{i_\nu}(x)\subseteq\mathcal R_e(x_{<i_0});
\]

\begin{itemize}
\tightlist
\item
  for every distribution on \(\mathcal W\),
\end{itemize}

\[
H(F)\le N\bigl[h(\theta)+\theta\log Q\bigr],
\qquad
h(t)=-t\log t-(1-t)\log(1-t).
\]

The window parameter \(q\) in the application only sets \(N=L-q+1\); the abstract theorem itself does not involve windows.

\subsection{Literals, truth, and complete origins}\label{literals-truth-and-complete-origins}

Fix one realized word \(x\) for the construction and counting arguments, and write \(z_i=Z_i(x)\). A \textbf{literal} is a pair \((i,a)\) with \(a\in\Omega\). A literal is true if \(a=z_i\). There are \(QN\) literals.

A \textbf{complete origin} consists of:

\begin{itemize}
\tightlist
\item
  one positional edge \(e=(s,i_1,\ldots,i_{r-1})\);
\item
  one assignment \((a_1,\ldots,a_{r-1})\) of letters to its future positions;
\item
  the condition
\end{itemize}

\[
(z_s,a_1,\ldots,a_{r-1})\notin\mathcal R_e(x_{<s}).
\]

The source argument is the one fixed true source letter. An origin is therefore a forbidden future set of \(r-1\) literals. It always contains a false literal: otherwise the realized word would violate its original relation.

A projection of an origin removes some true literals whose positions have been explicitly recorded. An origin becomes irrelevant when any of its required literals is excluded from a palette. An origin may also be discarded at the chronological retirement of an unrecorded position; this latter loss will be charged to the bad-edge budget.

\subsubsection{Pair multiplicity lemma}\label{pair-multiplicity-lemma}

For any two future literals \(v=(i,a)\), \(w=(j,c)\) at distinct positions, at most \(D=CQ^{r-3}\) complete origins contain both.

\textbf{Proof.} There are at most \(C\) original positional edges containing the two positions. In any origin whose future set contains both literals, its source is earlier than both and its source letter is already fixed to the true letter. The two prescribed future letters leave only \(r-3\) future letters unspecified. There are at most \(Q^{r-3}\) possibilities. Sum over the at most \(C\) edges. \(\square\)

This is a bound on distinct full origins over the entire execution for this word. It is not a bound on the number of historical projection shapes. The distinction is essential: an origin can appear successively at several ranks, but its complete label is unchanged.

One may define the universe of all complete origins at the outset for purposes of counting, even though the actual algorithm enumerates an origin only when its source is processed. Thus an origin selected just before synchronous palette deletions can be charged to those deletions even if its queue job has not yet been materialized.

\subsection{State and invariants}\label{state-and-invariants}

Each position \(i\) has a palette \(S_i\), initially \(\Omega\). Palettes only shrink, and the true letter is always retained.

A position is \textbf{recorded} when its true letter is put in the fingerprint. Its palette then becomes the singleton \(\{z_i\}\). A recorded true literal remains in that palette but is inactive in all graphs.

An unprocessed, unrecorded literal in its current palette is active. At the chronological retirement of a position, all its remaining active literals become inactive permanently. Inactivity and palette exclusion are different operations.

For each \(2\le j\le r-1\), maintain a simple \(j\)-uniform hypergraph \(J_j\) on active literals. Every physical hyperedge carries one or more complete-origin labels. All labels carried by a physical edge have that edge as their current projection.

Also maintain a finite queue of jobs. A job contains a projected set of literals, the full-origin labels it owns, and a frozen rank bound. Its actual remaining arity may fall below that bound when recorded literals are substituted or excluded literals are noticed. The bound never increases along an origin's history.

We maintain:

\textbf{I1. Truth preservation.} Every palette contains its true letter.

\textbf{I2. Unique ownership.} Each live origin has exactly one current owner: one physical edge or one queued job. Equal physical edges merge their label sets. Origins never branch.

\textbf{I3. Sound projection.} A live projection is obtained from its complete origin only by removing recorded true literals. No unrecorded true value is silently substituted.

\textbf{I4. No false resurrection.} A literal excluded from its palette never returns. An inactive literal never becomes active again.

\textbf{I5. Degree control.} Immediately before every physical insertion, all active degrees in each \(J_j\) are strictly less than \(T_j\). At the instant any active literal is removed, its degree in \(J_j\) is at most \(T_j\).

\textbf{I6. Source barriers.} Before the chronological retirement of a position, and before the next source is processed, every queue is empty.

All choices---edge order, literal order, witness order, and queue tie breaking---are fixed deterministically from the available state. The full realized word is used only to answer an explicitly posed equality query and to select the heavy or light actual source branch. Replay will eliminate the latter dependence.

\subsection{Primitive operations}\label{primitive-operations}

\subsubsection{Excluding a false literal}\label{excluding-a-false-literal}

Delete the literal from its palette and from the active graphs. Every active graph origin requiring it becomes irrelevant. A queued origin requiring it is likewise irrelevant, even if its job discovers this only later. Record the time of the actual palette exclusion for accounting; a later inspection of the job is not a second deletion event.

A negative equality query is one way to exclude a false literal. Unit propagation from a valid forbidden projection is another.

\subsubsection{Recording a true literal}\label{recording-a-true-literal}

For an active literal \(v=(i,z_i)\):

\begin{enumerate}
\def\labelenumi{\arabic{enumi}.}
\tightlist
\item
  first snapshot and take ownership of all physical edges incident to \(v\), in every rank, with their origin labels;
\item
  add \((i,z_i)\) to the fingerprint; exclude all other palette letters at position \(i\); deactivate \(v\), preserving its singleton palette;
\item
  for every snapshotted rank-two edge \(\{v,w\}\), exclude \(w\), before inserting any new physical edge;
\item
  for every other snapshotted edge \(B\), put \(B\setminus\{v\}\), with its origin labels, into a queue with frozen rank \(|B|-1\).
\end{enumerate}

The order can be implemented without ambiguity by first removing all snapshotted physical edges from the graph and giving their labels to their projected jobs or unit consequences. Exclusions can make some of those projected jobs irrelevant before they are later inspected.

Every rank-two neighbor is false: if it were true, an origin of the edge, together with the previously substituted true literals, would make the original forbidden assignment entirely true. Thus I1 is preserved.

Edges contain at most one literal from any position, because they are projections of positional edges. In particular, the other letters excluded at position \(i\) do not form second ends of the saved edges incident to \(v\).

\subsubsection{Processing a job}\label{processing-a-job}

If a required letter has been excluded from its palette, discard the job's origins. Otherwise substitute every recorded true literal in it. A letter at a recorded position that does not match the recorded value was already excluded and falls under the first case.

There cannot be zero remaining literals, because then a forbidden origin would be entirely true. With one remaining literal, exclude that false literal. With \(j\ge2\) remaining literals, insert their set into \(J_j\), or merge the origin labels if that physical edge already exists.

A merge does not increase any physical degree and is not a physical insertion.

\subsubsection{Degree-triggered queries}\label{degree-triggered-queries}

Suppose a new physical edge of rank \(j\) is inserted. Before insertion every degree is below \(T_j\), so only its ends can reach \(T_j\), and none can exceed it. If an end reaches \(T_j\), choose one such end \(v\) and query whether it is true.

A negative response excludes \(v\). A positive response records \(v\). In either case the newly inserted edge disappears immediately. Every other end that had just reached the threshold loses this same new edge and returns below threshold. At that moment other ranks have undergone only removals.

The only remaining operations that can increase degrees are later physical insertions, each handled by the same rule. This proves I5 inductively, once the scheduler below ensures that new insertions are processed one at a time.

\subsection{An explicit packet scheduler}\label{an-explicit-packet-scheduler}

A positive query that is triggered in rank \(j\) opens its rank-\(j\) packet \textbf{before} the snapshot and recording operation. Thus the packet includes synchronous exclusion of the other letters at the recorded position and all saved rank-two neighbors.

After recording, drain all pending jobs whose frozen rank is strictly less than \(j\). Each such job has actual arity at most its frozen rank, so any positive query it triggers has rank strictly less than \(j\). That query opens a nested packet of strictly smaller rank. On return, continue draining eligible jobs until none remains. Then close the rank-\(j\) packet.

Jobs of frozen rank at least \(j\) are deferred while this packet is open. A parent packet can process them later if their frozen rank is below the parent's bound. An external drain with bound \(+\infty\) processes all jobs.

This scheduler yields three facts used below:

\begin{enumerate}
\def\labelenumi{\arabic{enumi}.}
\tightlist
\item
  nested positive packets have strictly decreasing ranks;
\item
  two packets of the same rank have disjoint time intervals;
\item
  an origin that started as the projection of a saved rank-\(j\) edge can never be stranded in the rank-\(j\) packet's deferred queue: its initial projection has frozen rank \(j-1\), and every subsequent projection along that origin's chain has still smaller frozen rank.
\end{enumerate}

A job that came from an unrelated larger edge can have actual arity below \(j\) while retaining frozen rank at least \(j\); deferring it is harmless. The third fact concerns the selected own origins of the packet, whose frozen ranks are already below \(j\).

There is no chronological retirement inside a packet. Chronological retirement happens only after the external drain has finished.

Termination is finite. Only finitely many original forbidden assignments can be installed. Each origin has one chain, and every transfer created by recording strictly reduces its number of required literals. A job is processed once unless it is replaced by such a strict projection; it is otherwise discarded or attached to one physical edge. No origin is duplicated. There are therefore finitely many insertions, transfers, and query events.

\subsection{Processing sources}\label{processing-sources}

Process sources \(i=0,\ldots,N-1\) in order. Queues are empty at entry. The relations whose source is \(i\) are known from \(x_{<i}\).

If \(i\) is already recorded, assign \(A_i=\{z_i\}\), enumerate every forbidden future tuple with source letter \(z_i\) in the current future palettes, and externally drain all resulting jobs.

If \(i\) is not recorded, call a candidate \(a\in S_i\) \textbf{heavy} if at least \(T_{\rm s}\) outgoing positional edges admit a forbidden future tuple in the current future palettes when their first letter is \(a\). Otherwise it is light.

If the true letter is heavy:

\begin{itemize}
\tightlist
\item
  snapshot the current palettes and the full list of forbidden future tuples for the true source letter;
\item
  begin the source accounting interval before any recording or deletion;
\item
  record the source, assign \(A_i=\{z_i\}\), and enqueue the snapshotted list;
\item
  externally drain every queue.
\end{itemize}

If the true letter is light, assign \(A_i\) to be all light letters of \(S_i\). This nonempty set contains the true letter. This branch creates no new constraints, no new jobs, and no queries.

Finally, with all queues empty, retire the remaining active literals at position \(i\). Delete all physical edges incident to them, discarding their origins without substituting the unrecorded actual value. Do not change the assigned \(A_i\) afterward.

The future positions of any newly installed origin are later than its source. Thus a newly installed origin never contains an already retired position. At a chronological retirement, every earlier queued job has been processed; all still-live origins requiring that position are on its active incident edges and are discarded then. Consequently no later event refers to a retired position. This also proves that future palette changes never alter an already assigned alphabet.

\subsection{Physical insertion budget}\label{physical-insertion-budget}

Let \(I_j\) be the number of physical insertions in rank \(j\). Every inserted edge is eventually removed, because all its positions are ultimately retired or recorded.

Charge an insertion to the first removal of one of its active endpoints. Each literal is removed from active status once, and its degree at that event is at most \(T_j\). There are \(QN\) literals. Therefore

\[
I_j\le T_jQN. \tag{7.1}
\]

A physical edge with a given literal set cannot disappear and later be reinserted with that same literal set: its first disappearance removes at least one of those active literals permanently. This is also a direct reason why it is legitimate to speak of the unique insertion event of a surviving physical edge. Alternatively, charging by insertion occurrence gives the same estimate.

\subsection{Positive-query budgets: the central charge maps}\label{positive-query-budgets-the-central-charge-maps}

Let \(H_j\) be the number of positive rank-\(j\) queries, and let \(H_{\rm s}\) be the number of newly recorded heavy sources.

\subsubsection{Rank two}\label{rank-two}

At a positive rank-two query, the queried literal has exactly \(T_2\) distinct physical neighbors. All are active and present in their palettes at that instant. Recording removes these \(T_2\) distinct false neighbors before any new insertion.

Each literal can be excluded from its palette at most once. Hence

\[
H_2T_2\le QN. \tag{8.1}
\]

Neighbors may belong to the same position but be different letters; they are still different literals, which is exactly what the \(QN\) budget counts.

\subsubsection{\texorpdfstring{Rank \(j\ge3\): selection}{Rank j\textbackslash ge3: selection}}\label{rank-jge3-selection}

At a positive rank-\(j\) query with true pivot \(v\), its snapshot has exactly \(T_j\) incident physical edges. Select one complete-origin label from each. These labels are distinct by unique ownership.

Immediately after substituting \(v\), every selected origin has a projection of rank \(j-1\). Such an origin is never selected again by a different rank-\(j\) packet, because its single chain can never return to rank \(j\). Thus there are exactly \(H_jT_j\) distinct rank-\(j\) selected-origin instances over all such packets.

\subsubsection{\texorpdfstring{Rank \(j\ge3\): exhaustive exit alternatives}{Rank j\textbackslash ge3: exhaustive exit alternatives}}\label{rank-jge3-exhaustive-exit-alternatives}

For each selected origin, at the closure of its packet exactly one of the following accounting alternatives applies:

\begin{itemize}
\tightlist
\item
  it has lost a required palette letter during this packet;
\item
  it is carried by an active physical edge of rank \(l<j\).
\end{itemize}

Here is why this list is exhaustive. No chronological retirement occurs within the packet. The origin started alive and its actual projection has rank \(j-1\). Every subsequent job on its chain has frozen rank below \(j\), so cannot remain in the deferred queue when the packet closes. If a recorded end is substituted, the rank decreases further. A zero-end origin is impossible by soundness. A one-end origin causes a palette exclusion. If it is neither irrelevant nor still queued, unique ownership places it on an active physical edge, necessarily of smaller rank.

\subsubsection{Charging actual palette losses}\label{charging-actual-palette-losses}

For the first alternative, charge the selected origin to its first required palette exclusion during the packet. That excluded literal \(w\) and the packet pivot \(v\) both belong to the origin's full future assignment. They are distinct: the true pivot is recorded and retained, not excluded.

For one fixed packet and one fixed exclusion event \(w\), the pair lemma bounds these selected origins by \(D\). Packets of the same rank are time-disjoint, so an actual exclusion event belongs to at most one rank-\(j\) packet. There are at most \(QN\) exclusion events overall. The total loss charge is therefore at most \(DQN\).

This includes exclusions at packet opening. An origin whose queued projection is inspected later is charged to the earlier actual exclusion event within the same packet. Nothing is charged to a merely administrative late inspection.

\subsubsection{Charging surviving lower edges}\label{charging-surviving-lower-edges}

For the second alternative, charge the selected origin to the unique insertion event of its final active physical edge \(B\) of rank \(l<j\).

Fix any pair of literals in \(B\). Every complete origin charged to this edge contains that pair. Across all rank-\(j\) packets, each complete origin is selected at most once at rank \(j\), and at most \(D\) complete origins contain the pair. The physical insertion consequently pays for at most \(D\) selected origins, even if some origins merge into an edge created before their packets began.

This is the precise resolution of the ``old lower edge receives new projections'' concern. The budget counts its full-origin labels over time, not its creation time relative to the paying packet.

Combining both charge classes gives

\[
H_jT_j\le D\left(QN+\sum_{l=2}^{j-1}I_l\right),
\qquad 3\le j\le r-1. \tag{8.2}
\]

\subsection{Heavy-source budget}\label{heavy-source-budget}

At a newly recorded heavy source choose \(T_{\rm s}\) different outgoing positional edges witnessing heaviness, and one forbidden future tuple from each, in the palettes snapshotted before recording.

Their future position sets are pairwise disjoint. These tuples are alive at the start of the source accounting interval. At the end of the external drain, a selected origin either has lost a palette letter during this interval or is on an active physical edge. There are no queued jobs, no chronological retirements during the interval, and no completely true forbidden tuple.

An actual excluded literal belongs to at most one selected future tuple for this source, by future-position disjointness. Source intervals are successive and disjoint. Thus palette-loss charges total at most \(QN\).

For surviving-edge charges, any fixed pair of the active edge belongs to every paying complete origin; a source is processed only once and one origin was chosen for each selected outgoing edge. The pair lemma bounds charges to a physical insertion by \(D\), over all source intervals.

Therefore

\[
H_{\rm s}T_{\rm s}\le QN+D\sum_{j=2}^{r-1}I_j. \tag{9.1}
\]

It is legitimate for a deletion or insertion to participate in the separate upper bounds for different ranks and for heavy sources. These are different estimates for different record counts; each estimate has its own valid charge map. The proof never asserts a single shared global budget across all these independently bounded classes.

\subsection{Fingerprint size}\label{fingerprint-size}

A position is recorded only once. Every recording is either a new heavy-source recording or a positive query recording. Thus

\[
|F|=H_{\rm s}+\sum_{j=2}^{r-1}H_j.
\]

Using (7.1), (8.1), (8.2), and (9.1) gives

\[
H_{\rm s}\le QN\frac{D\sum_{j=2}^{r-1}T_j+1}{T_{\rm s}},
\]

\[
H_2\le QN/T_2,
\]

\[
H_j\le DQN\frac{1+\sum_{l=2}^{j-1}T_l}{T_j}
\quad (j\ge3).
\]

Their sum is exactly \(\theta N\), with the stated constant. No hidden dependence on \(N\), the relations, or the realized word appears.

\subsection{Bad-edge budget}\label{bad-edge-budget}

Call an edge bad for this word if its assigned full Cartesian product fails to be contained in the relation frozen at its actual source prefix.

\subsubsection{Light sources}\label{light-sources}

For a light source, each of its at most \(Q\) assigned source candidates has fewer than \(T_{\rm s}\) outgoing edges with a forbidden future tuple in the then-current future palettes. Later future alphabets are subsets of those palettes. Consequently the source has fewer than \(QT_{\rm s}\) bad outgoing edges.

The total light-source contribution is at most \(QT_{\rm s}N\).

\subsubsection{Recorded sources}\label{recorded-sources}

For a bad edge with recorded source, choose one forbidden tuple from the final assigned alphabets. Its source letter is the actual recorded letter. Its future letters were all present when its source was processed, because palettes shrink and assigned alphabets are never enlarged after retirement. Hence this full origin was enumerated or was already rendered irrelevant only by a palette event that would also exclude a final alphabet letter.

The latter cannot happen: every chosen future letter belongs to its assigned final alphabet, so none was ever excluded from its palette. Recording one of its future positions can only substitute the matching true letter; recording a different letter would exclude the chosen one, again contradicting final membership.

Follow the origin's unique chain. It cannot end in a unit exclusion or an empty projection, and it cannot disappear through a false-letter exclusion. At the end of the complete execution no active literal remains. Because every chronological retirement has an empty queue, the origin must first be discarded on an active physical edge at a chronological retirement.

Charge this bad positional edge to that physical edge's insertion event, using the selected complete origin as the invoice. Different bad positional edges have different complete origins. Every origin charged to a fixed physical edge contains any chosen pair of its literals. The pair lemma therefore bounds its charges by \(D\).

The recorded-source contribution is at most \(D\sum_j I_j\). Together with (7.1),

\[
E_{\rm bad}
\le QT_{\rm s}N+D\sum_{j=2}^{r-1}I_j
\le Q\left(T_{\rm s}+D\sum_{j=2}^{r-1}T_j\right)N.
\]

This bound explicitly pays for chronological deletion without substituting an unrecorded actual value. Omitting that payment would leave a real gap; it is included in the bound here.

\subsection{Exact replay from the final fingerprint}\label{exact-replay-from-the-final-fingerprint}

Fix a realized pair \((x_{<i},F(x))\). We show that the entire execution through assignment of \(A_i\) can be reconstructed from this pair.

At any equality query on an active literal \((k,a)\), the answer is positive if and only if \((k,a)\in F(x)\):

\begin{itemize}
\tightlist
\item
  every positive answer immediately records the pair;
\item
  every recorded pair is true;
\item
  a true literal cannot be excluded or change its value later.
\end{itemize}

The fingerprint is used only to answer the queries actually produced by the deterministic algorithm. It is not used to insert all future recorded letters into the state in advance. Absence of an unqueried pair says nothing by itself about that letter.

Process earlier sources by induction. Their relations are known because their source prefixes are restrictions of \(x_{<i}\). Each queue action and literal query is then determined by the reconstructed state and the fingerprint oracle.

The only remaining branch is a source \(s\le i\) that has not already been recorded when its turn arrives. At that point all queues are empty.

\begin{itemize}
\tightlist
\item
  If \(s\notin\operatorname{dom}F\), its actual source letter cannot be heavy, because that would record it. Its actual letter is light. The algorithm computes \(A_s\) as all light letters from the known state and creates no events.
\item
  If \(s\in\operatorname{dom}F\), its true letter is given by its fingerprint pair. If that letter were light, the light branch would create no events; the source would immediately retire, after which it could never be recorded. This contradicts membership in the final fingerprint. Therefore the actual letter is heavy and the heavy branch is determined.
\end{itemize}

An already recorded source uses its known recorded value. This completes the induction. The algorithm determines \(A_i\) from \((x_{<i},F)\) exactly as claimed.

Consequently, for any probability distribution on \(\mathcal W\), the assigned \(A_i\) contains the entire conditional support

\[
\operatorname{supp}(Z_i\mid x_{<i},F),
\]

on every conditioning atom of positive probability. This is stronger than merely containing the one actual letter on one execution.

\subsection{Fingerprint entropy}\label{fingerprint-entropy}

There are at most

\[
\sum_{a=0}^{\lfloor\theta N\rfloor}\binom NaQ^a
\]

possible fingerprints. Since \(\theta\le1/2\),

\[
\sum_{a\le\lfloor\theta N\rfloor}\binom Na
\le \exp\{Nh(\theta)\},
\qquad
Q^a\le Q^{\theta N}.
\]

Thus

\[
H(F)\le\log|\operatorname{range}F|
\le N[h(\theta)+\theta\log Q].
\]

The map \(F\) is deterministic, but this count gives the same bound for every distribution on \(\mathcal W\). The proof of the finite theorem is complete.


\end{document}
